% Options for packages loaded elsewhere
% Options for packages loaded elsewhere
\PassOptionsToPackage{unicode}{hyperref}
\PassOptionsToPackage{hyphens}{url}
\PassOptionsToPackage{dvipsnames,svgnames,x11names}{xcolor}
%
\documentclass[
  portuguese,
  letterpaper,
  openany,
  oneside,
  12pt,
  a4paper]{book}
\usepackage{xcolor}
\usepackage{amsmath,amssymb}
\setcounter{secnumdepth}{5}
\usepackage{iftex}
\ifPDFTeX
  \usepackage[T1]{fontenc}
  \usepackage[utf8]{inputenc}
  \usepackage{textcomp} % provide euro and other symbols
\else % if luatex or xetex
  \usepackage{unicode-math} % this also loads fontspec
  \defaultfontfeatures{Scale=MatchLowercase}
  \defaultfontfeatures[\rmfamily]{Ligatures=TeX,Scale=1}
\fi
\usepackage{lmodern}
\ifPDFTeX\else
  % xetex/luatex font selection
\fi
% Use upquote if available, for straight quotes in verbatim environments
\IfFileExists{upquote.sty}{\usepackage{upquote}}{}
\IfFileExists{microtype.sty}{% use microtype if available
  \usepackage[]{microtype}
  \UseMicrotypeSet[protrusion]{basicmath} % disable protrusion for tt fonts
}{}
\makeatletter
\@ifundefined{KOMAClassName}{% if non-KOMA class
  \IfFileExists{parskip.sty}{%
    \usepackage{parskip}
  }{% else
    \setlength{\parindent}{0pt}
    \setlength{\parskip}{6pt plus 2pt minus 1pt}}
}{% if KOMA class
  \KOMAoptions{parskip=half}}
\makeatother
% Make \paragraph and \subparagraph free-standing
\makeatletter
\ifx\paragraph\undefined\else
  \let\oldparagraph\paragraph
  \renewcommand{\paragraph}{
    \@ifstar
      \xxxParagraphStar
      \xxxParagraphNoStar
  }
  \newcommand{\xxxParagraphStar}[1]{\oldparagraph*{#1}\mbox{}}
  \newcommand{\xxxParagraphNoStar}[1]{\oldparagraph{#1}\mbox{}}
\fi
\ifx\subparagraph\undefined\else
  \let\oldsubparagraph\subparagraph
  \renewcommand{\subparagraph}{
    \@ifstar
      \xxxSubParagraphStar
      \xxxSubParagraphNoStar
  }
  \newcommand{\xxxSubParagraphStar}[1]{\oldsubparagraph*{#1}\mbox{}}
  \newcommand{\xxxSubParagraphNoStar}[1]{\oldsubparagraph{#1}\mbox{}}
\fi
\makeatother

\usepackage{color}
\usepackage{fancyvrb}
\newcommand{\VerbBar}{|}
\newcommand{\VERB}{\Verb[commandchars=\\\{\}]}
\DefineVerbatimEnvironment{Highlighting}{Verbatim}{commandchars=\\\{\}}
% Add ',fontsize=\small' for more characters per line
\usepackage{framed}
\definecolor{shadecolor}{RGB}{241,243,245}
\newenvironment{Shaded}{\begin{snugshade}}{\end{snugshade}}
\newcommand{\AlertTok}[1]{\textcolor[rgb]{0.68,0.00,0.00}{#1}}
\newcommand{\AnnotationTok}[1]{\textcolor[rgb]{0.37,0.37,0.37}{#1}}
\newcommand{\AttributeTok}[1]{\textcolor[rgb]{0.40,0.45,0.13}{#1}}
\newcommand{\BaseNTok}[1]{\textcolor[rgb]{0.68,0.00,0.00}{#1}}
\newcommand{\BuiltInTok}[1]{\textcolor[rgb]{0.00,0.23,0.31}{#1}}
\newcommand{\CharTok}[1]{\textcolor[rgb]{0.13,0.47,0.30}{#1}}
\newcommand{\CommentTok}[1]{\textcolor[rgb]{0.37,0.37,0.37}{#1}}
\newcommand{\CommentVarTok}[1]{\textcolor[rgb]{0.37,0.37,0.37}{\textit{#1}}}
\newcommand{\ConstantTok}[1]{\textcolor[rgb]{0.56,0.35,0.01}{#1}}
\newcommand{\ControlFlowTok}[1]{\textcolor[rgb]{0.00,0.23,0.31}{\textbf{#1}}}
\newcommand{\DataTypeTok}[1]{\textcolor[rgb]{0.68,0.00,0.00}{#1}}
\newcommand{\DecValTok}[1]{\textcolor[rgb]{0.68,0.00,0.00}{#1}}
\newcommand{\DocumentationTok}[1]{\textcolor[rgb]{0.37,0.37,0.37}{\textit{#1}}}
\newcommand{\ErrorTok}[1]{\textcolor[rgb]{0.68,0.00,0.00}{#1}}
\newcommand{\ExtensionTok}[1]{\textcolor[rgb]{0.00,0.23,0.31}{#1}}
\newcommand{\FloatTok}[1]{\textcolor[rgb]{0.68,0.00,0.00}{#1}}
\newcommand{\FunctionTok}[1]{\textcolor[rgb]{0.28,0.35,0.67}{#1}}
\newcommand{\ImportTok}[1]{\textcolor[rgb]{0.00,0.46,0.62}{#1}}
\newcommand{\InformationTok}[1]{\textcolor[rgb]{0.37,0.37,0.37}{#1}}
\newcommand{\KeywordTok}[1]{\textcolor[rgb]{0.00,0.23,0.31}{\textbf{#1}}}
\newcommand{\NormalTok}[1]{\textcolor[rgb]{0.00,0.23,0.31}{#1}}
\newcommand{\OperatorTok}[1]{\textcolor[rgb]{0.37,0.37,0.37}{#1}}
\newcommand{\OtherTok}[1]{\textcolor[rgb]{0.00,0.23,0.31}{#1}}
\newcommand{\PreprocessorTok}[1]{\textcolor[rgb]{0.68,0.00,0.00}{#1}}
\newcommand{\RegionMarkerTok}[1]{\textcolor[rgb]{0.00,0.23,0.31}{#1}}
\newcommand{\SpecialCharTok}[1]{\textcolor[rgb]{0.37,0.37,0.37}{#1}}
\newcommand{\SpecialStringTok}[1]{\textcolor[rgb]{0.13,0.47,0.30}{#1}}
\newcommand{\StringTok}[1]{\textcolor[rgb]{0.13,0.47,0.30}{#1}}
\newcommand{\VariableTok}[1]{\textcolor[rgb]{0.07,0.07,0.07}{#1}}
\newcommand{\VerbatimStringTok}[1]{\textcolor[rgb]{0.13,0.47,0.30}{#1}}
\newcommand{\WarningTok}[1]{\textcolor[rgb]{0.37,0.37,0.37}{\textit{#1}}}

\usepackage{longtable,booktabs,array}
\usepackage{calc} % for calculating minipage widths
% Correct order of tables after \paragraph or \subparagraph
\usepackage{etoolbox}
\makeatletter
\patchcmd\longtable{\par}{\if@noskipsec\mbox{}\fi\par}{}{}
\makeatother
% Allow footnotes in longtable head/foot
\IfFileExists{footnotehyper.sty}{\usepackage{footnotehyper}}{\usepackage{footnote}}
\makesavenoteenv{longtable}
\usepackage{graphicx}
\makeatletter
\newsavebox\pandoc@box
\newcommand*\pandocbounded[1]{% scales image to fit in text height/width
  \sbox\pandoc@box{#1}%
  \Gscale@div\@tempa{\textheight}{\dimexpr\ht\pandoc@box+\dp\pandoc@box\relax}%
  \Gscale@div\@tempb{\linewidth}{\wd\pandoc@box}%
  \ifdim\@tempb\p@<\@tempa\p@\let\@tempa\@tempb\fi% select the smaller of both
  \ifdim\@tempa\p@<\p@\scalebox{\@tempa}{\usebox\pandoc@box}%
  \else\usebox{\pandoc@box}%
  \fi%
}
% Set default figure placement to htbp
\def\fps@figure{htbp}
\makeatother



\ifLuaTeX
\usepackage[bidi=basic]{babel}
\else
\usepackage[bidi=default]{babel}
\fi
% get rid of language-specific shorthands (see #6817):
\let\LanguageShortHands\languageshorthands
\def\languageshorthands#1{}


\setlength{\emergencystretch}{3em} % prevent overfull lines

\providecommand{\tightlist}{%
  \setlength{\itemsep}{0pt}\setlength{\parskip}{0pt}}



 
\usepackage[style=abnt,backend=biber,style=abnt,citestyle=abnt]{biblatex}
\addbibresource{/Users/fz/VSCode/pdi-vc/references.bib}


% preamble.tex — Configurações LaTeX para o livro PDI+VC
\usepackage[T1]{fontenc}
\usepackage[utf8]{inputenc}
\usepackage{lmodern}
\usepackage{microtype}
\usepackage{amsmath,amssymb}
\usepackage{booktabs}
\usepackage{longtable}
\usepackage{array}
\usepackage{graphicx}
\usepackage{float}
\usepackage{subcaption}
\usepackage{xcolor}
\usepackage{listings}
\usepackage{fancyvrb}
\usepackage{hyperref}

% Configuração de código-fonte
\lstset{
  basicstyle=\ttfamily\small,
  breaklines=true,
  breakatwhitespace=true,
  showstringspaces=false,
  frame=single,
  numbers=left,
  numberstyle=\tiny\color{gray},
  keywordstyle=\color{blue}\bfseries,
  commentstyle=\color{green!50!black},
  stringstyle=\color{red!70!black},
}

\lstdefinestyle{python}{
  language=Python,
  morekeywords={import,from,as,def,class,return,yield,with,lambda},
}

\lstdefinestyle{cpp}{
  language=C++,
  morekeywords={cout,cin,endl,Mat,Vec3b,uchar,using,namespace,cv},
}

% Cores PDI+VC
\definecolor{pdi-blue}{RGB}{21,101,192}
\definecolor{pdi-green}{RGB}{46,125,50}
\makeatletter
\@ifpackageloaded{tcolorbox}{}{\usepackage[skins,breakable]{tcolorbox}}
\@ifpackageloaded{fontawesome5}{}{\usepackage{fontawesome5}}
\definecolor{quarto-callout-color}{HTML}{909090}
\definecolor{quarto-callout-note-color}{HTML}{0758E5}
\definecolor{quarto-callout-important-color}{HTML}{CC1914}
\definecolor{quarto-callout-warning-color}{HTML}{EB9113}
\definecolor{quarto-callout-tip-color}{HTML}{00A047}
\definecolor{quarto-callout-caution-color}{HTML}{FC5300}
\definecolor{quarto-callout-color-frame}{HTML}{acacac}
\definecolor{quarto-callout-note-color-frame}{HTML}{4582ec}
\definecolor{quarto-callout-important-color-frame}{HTML}{d9534f}
\definecolor{quarto-callout-warning-color-frame}{HTML}{f0ad4e}
\definecolor{quarto-callout-tip-color-frame}{HTML}{02b875}
\definecolor{quarto-callout-caution-color-frame}{HTML}{fd7e14}
\makeatother
\makeatletter
\@ifpackageloaded{bookmark}{}{\usepackage{bookmark}}
\makeatother
\makeatletter
\@ifpackageloaded{caption}{}{\usepackage{caption}}
\AtBeginDocument{%
\ifdefined\contentsname
  \renewcommand*\contentsname{Índice}
\else
  \newcommand\contentsname{Índice}
\fi
\ifdefined\listfigurename
  \renewcommand*\listfigurename{Lista de Figuras}
\else
  \newcommand\listfigurename{Lista de Figuras}
\fi
\ifdefined\listtablename
  \renewcommand*\listtablename{Lista de Tabelas}
\else
  \newcommand\listtablename{Lista de Tabelas}
\fi
\ifdefined\figurename
  \renewcommand*\figurename{Figura}
\else
  \newcommand\figurename{Figura}
\fi
\ifdefined\tablename
  \renewcommand*\tablename{Tabela}
\else
  \newcommand\tablename{Tabela}
\fi
}
\@ifpackageloaded{float}{}{\usepackage{float}}
\floatstyle{ruled}
\@ifundefined{c@chapter}{\newfloat{codelisting}{h}{lop}}{\newfloat{codelisting}{h}{lop}[chapter]}
\floatname{codelisting}{Listagem}
\newcommand*\listoflistings{\listof{codelisting}{Lista de Listagens}}
\makeatother
\makeatletter
\makeatother
\makeatletter
\@ifpackageloaded{caption}{}{\usepackage{caption}}
\@ifpackageloaded{subcaption}{}{\usepackage{subcaption}}
\makeatother
\usepackage{bookmark}
\IfFileExists{xurl.sty}{\usepackage{xurl}}{} % add URL line breaks if available
\urlstyle{same}
\hypersetup{
  pdftitle={Processamento Digital de Imagens e Visão Computacional},
  pdfauthor={Francisco de Assis Zampirolli},
  pdflang={pt},
  colorlinks=true,
  linkcolor={blue},
  filecolor={Maroon},
  citecolor={Blue},
  urlcolor={blue},
  pdfcreator={LaTeX via pandoc}}


\title{Processamento Digital de Imagens e Visão Computacional}
\usepackage{etoolbox}
\makeatletter
\providecommand{\subtitle}[1]{% add subtitle to \maketitle
  \apptocmd{\@title}{\par {\large #1 \par}}{}{}
}
\makeatother
\subtitle{Abordagem Prática com Python}
\author{Francisco de Assis Zampirolli}
\date{2026-05-06}
\begin{document}
\frontmatter
\maketitle

\renewcommand*\contentsname{Índice}
{
\hypersetup{linkcolor=}
\setcounter{tocdepth}{2}
\tableofcontents
}

\mainmatter
\bookmarksetup{startatroot}

\chapter*{PDI e VC --- Python}\label{pdi-e-vc-python}
\addcontentsline{toc}{chapter}{PDI e VC --- Python}

\markboth{PDI e VC --- Python}{PDI e VC --- Python}

Bem-vindo ao livro de PDI e Visão Computacional --- versão Python /
Português.

\subsection*{Organização}\label{organizauxe7uxe3o}
\addcontentsline{toc}{subsection}{Organização}

\begin{itemize}
\tightlist
\item
  \textbf{Parte I --- Fundamentos de PDI} --- Representação,
  histogramas, filtragem, morfologia
\item
  \textbf{Parte II --- Visão Computacional} --- Segmentação,
  descritores, detecção, deep learning
\end{itemize}

\begin{center}\rule{0.5\linewidth}{0.5pt}\end{center}

\bookmarksetup{startatroot}

\chapter*{Prefácio}\label{prefuxe1cio}
\addcontentsline{toc}{chapter}{Prefácio}

\markboth{Prefácio}{Prefácio}

Este livro é um \textbf{projeto em construção} sobre
\textbf{Processamento Digital de Imagens (PDI) e Visão Computacional
(VC)}, concebido como material didático interativo para cursos de
graduação e pós-graduação em Computação, Engenharias e áreas afins.

\begin{quote}
📘 \textbf{Destaque:} Este livro segue a metodologia de
\textcite{zampirolli2025morph} --- uma versão estendida de
\textcite{zampirolli2024practical}, que recebeu o \textbf{prêmio de
melhor artigo} na trilha de ``Recursos e Ambientes Educacionais'' na
\textbf{EduComp 2024} --- com a biblioteca \texttt{morph.py},
desenvolvida pelo autor.
\end{quote}

\section*{Um livro vivo}\label{um-livro-vivo}
\addcontentsline{toc}{section}{Um livro vivo}

\markright{Um livro vivo}

Este não é um livro estático. O conteúdo está em constante evolução:
exemplos são refinados, novas seções são adicionadas e abordagens
pedagógicas são aprimoradas com base no feedback de alunos e
professores. Por isso, tratamos este material como um \emph{projeto
inicial} --- uma base sólida que continuará crescendo.

\section*{Uso de ferramentas de Inteligência
Artificial}\label{uso-de-ferramentas-de-inteliguxeancia-artificial}
\addcontentsline{toc}{section}{Uso de ferramentas de Inteligência
Artificial}

\markright{Uso de ferramentas de Inteligência Artificial}

Este livro foi produzido com o apoio de \textbf{ferramentas de
Inteligência Artificial} (IA) para aprimorar a qualidade do texto e do
código. A versão inicial dos capítulos passou por um processo de revisão
e correção utilizando múltiplos assistentes de IA, incluindo
\textbf{DeepSeek}, \textbf{Gemini} e \textbf{Claude}. Por meio de
\emph{prompts} iterativos, essas ferramentas ajudaram a:

\begin{itemize}
\tightlist
\item
  Refinar a clareza e a precisão das explicações;
\item
  Detectar e corrigir erros de sintaxe, lógica e digitação nos exemplos
  de código;
\item
  Sugerir melhorias na organização e na fluência do texto.
\end{itemize}

Além disso, \textbf{algumas imagens e ilustrações presentes no livro
foram geradas pelo Gemini}, sempre com o objetivo de facilitar a
compreensão de conceitos abstratos da área.

\begin{quote}
💡 \textbf{Transparência:} O uso de IA não substitui a revisão final do
autor, nem a validação pelos colaboradores humanos. O material é
constantemente testado e ajustado para garantir correção técnica e
adequação didática. Acreditamos que a parceria entre autor e ferramentas
de IA enriquece o processo de criação, desde que usada de forma ética e
criteriosa.
\end{quote}

\section*{Como este livro é
produzido}\label{como-este-livro-uxe9-produzido}
\addcontentsline{toc}{section}{Como este livro é produzido}

\markright{Como este livro é produzido}

Todo o conteúdo é escrito em \textbf{Quarto}, um sistema de publicação
científica e técnica que permite renderizar o mesmo código-fonte para
múltiplos formatos:

\begin{longtable}[]{@{}
  >{\raggedright\arraybackslash}p{(\linewidth - 2\tabcolsep) * \real{0.4500}}
  >{\raggedright\arraybackslash}p{(\linewidth - 2\tabcolsep) * \real{0.5500}}@{}}
\toprule\noalign{}
\begin{minipage}[b]{\linewidth}\raggedright
Formato
\end{minipage} & \begin{minipage}[b]{\linewidth}\raggedright
Descrição
\end{minipage} \\
\midrule\noalign{}
\endhead
\bottomrule\noalign{}
\endlastfoot
\textbf{HTML} & Versão web completa, com navegação interativa,
disponível em
\href{https://fzampirolli.github.io/pdi-vc/}{fzampirolli.github.io/pdi-vc} \\
\textbf{PDF} & Versão para impressão ou leitura offline \\
\textbf{Notebooks para alunos} & Versão \texttt{.ipynb} processada por
um filtro personalizado que resolve citações, numera figuras, tabelas e
equações, e gera referências formatadas no estilo ABNT --- ideal para
uso no Jupyter e Google Colab \\
\end{longtable}

O filtro de pós-processamento (\texttt{gerar\_notebooks\_alunos.py})
prepara os notebooks para distribuição aos alunos, garantindo que
funcionem perfeitamente em ambientes interativos mesmo sem depender do
Quarto instalado.

\section*{Sobre as linguagens de programação e
idiomas}\label{sobre-as-linguagens-de-programauxe7uxe3o-e-idiomas}
\addcontentsline{toc}{section}{Sobre as linguagens de programação e
idiomas}

\markright{Sobre as linguagens de programação e idiomas}

O \textbf{núcleo original} deste livro é escrito em \textbf{Português}
com exemplos em \textbf{Python}. Esta combinação (\texttt{pt\ +\ py}) é
a versão de referência, mantida e revisada diretamente pelos autor.

Para ampliar o alcance do material, versões em outras combinações de
idioma e linguagem de programação são geradas com o auxílio de
\textbf{APIs de Inteligência Artificial} (modelos de tradução e
conversão de código). Essas versões passam por um pipeline automatizado
que:

\begin{enumerate}
\def\labelenumi{\arabic{enumi}.}
\tightlist
\item
  \textbf{Traduz} o texto explicativo para o idioma de destino
\item
  \textbf{Converte} os exemplos de código Python para a linguagem alvo
  (ex: C++, R, Julia)
\item
  \textbf{Aplica} o mesmo pós-processamento de citações, figuras,
  tabelas e equações
\end{enumerate}

\subsection*{Exemplos de uso}\label{exemplos-de-uso}
\addcontentsline{toc}{subsection}{Exemplos de uso}

O projeto utiliza um \texttt{Makefile} com atalhos para desenvolvimento.
Veja alguns comandos comuns:

\begin{Shaded}
\begin{Highlighting}[]
\CommentTok{\# Versão de referência: Português + Python, HTML + PDF}
\FunctionTok{make}\NormalTok{ build{-}all}

\CommentTok{\# Apenas HTML (atualiza automaticamente ao salvar arquivos)}
\FunctionTok{make}\NormalTok{ html}

\CommentTok{\# Português + C++, apenas PDF}
\FunctionTok{make}\NormalTok{ build{-}pdf LANGS=cpp LOCALES=pt}

\CommentTok{\# Inglês + Python, apenas HTML}
\FunctionTok{make}\NormalTok{ build LANGS=py LOCALES=en}

\CommentTok{\# Inglês + C++, HTML + PDF}
\FunctionTok{make}\NormalTok{ build{-}all LANGS=cpp LOCALES=en}

\CommentTok{\# Todas as combinações disponíveis (py, cpp, java, c × pt, en, fr, es, it)}
\FunctionTok{make}\NormalTok{ full}

\CommentTok{\# Gerar notebooks para alunos (distribuição Jupyter/Colab)}
\FunctionTok{make}\NormalTok{ alunos}

\CommentTok{\# Publicar no GitHub Pages}
\FunctionTok{make}\NormalTok{ publish}
\end{Highlighting}
\end{Shaded}

\begin{quote}
💡 \textbf{Dica:} Você pode combinar qualquer idioma com qualquer
linguagem. Os idiomas disponíveis são: \texttt{pt} (Português),
\texttt{en} (Inglês), \texttt{fr} (Francês), \texttt{es} (Espanhol),
\texttt{it} (Italiano). As linguagens disponíveis são: \texttt{py}
(Python), \texttt{cpp} (C++), \texttt{java} (Java), \texttt{c} (C).
\end{quote}

\begin{quote}
⚠️ \textbf{Nota importante:} As versões geradas por IA são um
\textbf{ponto de partida} e podem conter imprecisões. Elas são
disponibilizadas como material complementar e devem ser revisadas antes
do uso em sala de aula. Contribuições da comunidade para refinar essas
versões são muito bem-vindas!
\end{quote}

\section*{Código aberto}\label{cuxf3digo-aberto}
\addcontentsline{toc}{section}{Código aberto}

\markright{Código aberto}

Este livro é um projeto de código aberto. Todo o conteúdo --- texto,
código-fonte, imagens e scripts de processamento --- está disponível
publicamente em:

👉
\textbf{\href{https://github.com/fzampirolli/pdi-vc/}{github.com/fzampirolli/pdi-vc}}

Você pode:

\begin{itemize}
\tightlist
\item
  📖 \textbf{Ler e estudar} gratuitamente
\item
  🐛 \textbf{Reportar erros} ou sugerir melhorias abrindo uma
  \emph{issue}
\item
  🤝 \textbf{Contribuir} com correções, novos exemplos ou traduções via
  \emph{pull requests}
\item
  🔧 \textbf{Adaptar} o material para suas próprias turmas
\end{itemize}

\section*{Como usar os notebooks}\label{como-usar-os-notebooks}
\addcontentsline{toc}{section}{Como usar os notebooks}

\markright{Como usar os notebooks}

Cada capítulo do livro está disponível como um \emph{notebook}
interativo. Se você nunca usou um notebook Jupyter ou Google Colab, veja
a seção \textbf{``Antes de começar: Notebooks em Python''} a seguir para
um guia rápido.

\section*{Agradecimentos}\label{agradecimentos}
\addcontentsline{toc}{section}{Agradecimentos}

\markright{Agradecimentos}

Agradecemos aos alunos, colegas e colaboradores que testam, questionam e
contribuem com este projeto. Este livro é feito para vocês e por vocês.

\begin{center}\rule{0.5\linewidth}{0.5pt}\end{center}

\section*{Antes de começar: Notebooks em
Python}\label{antes-de-comeuxe7ar-notebooks-em-python}
\addcontentsline{toc}{section}{Antes de começar: Notebooks em Python}

\markright{Antes de começar: Notebooks em Python}

Você já ouviu falar em \textbf{\emph{Literate Programming}} (Programação
Literária)? O conceito foi criado por Donald Knuth em 1984
\autocite{knuth_literate_1984} --- o mesmo que desenvolveu o sistema
\textbf{TeX} para tipografia digital. Knuth iniciou esse projeto no
final da década de 70 após ficar insatisfeito com a qualidade gráfica
das novas edições de sua série clássica, \emph{The Art of Computer
Programming} \autocite{knuth_art_1997}, buscando uma forma de unir a
precisão matemática à beleza estética dos livros antigos. O conceito de
Programação Literária não propõe apenas ``combinar código e
documentação'', mas sim inverter a lógica: o programa deve ser escrito
para ser lido por humanos (como um ensaio), e as partes de código são
extraídas dele para a máquina. É um detalhe sutil, mas que reforça o
termo ``literária''.

O conteúdo deste livro é formado por \emph{Notebook}: um documento que
intercala textos explicativos --- as \textbf{células de texto} no
formato
\href{https://colab.research.google.com/notebooks/markdown_guide.ipynb}{Markdown}
--- com códigos de programas em \href{https://www.python.org/}{Python}
--- as \textbf{células de código}. Mas aceita também outras linguagens
de programação, mas o processo de execução é mais complicado.

É fácil distinguir as células de código: elas são precedidas por
\texttt{{[}\ {]}}. Para executar uma célula de código, basta
selecioná-la e pressionar Shift + Enter. Antes de começar, algumas
observações importantes sobre o ambiente:

\begin{itemize}
\tightlist
\item
  \textbf{No computador:} no Colab ou Jupyter, você também pode executar
  uma célula clicando no botão de ▶️ (\emph{play}) que aparece ao passar
  o mouse sobre os colchetes \texttt{{[}\ {]}}.
\item
  \textbf{No celular:} no Colab, o atalho Shift + Enter pode não
  funcionar. Nesse caso, toque no botão de ▶️ (\emph{play}).
\item
  \textbf{Execução local:} você pode rodar este notebook no seu
  computador instalando o \href{https://jupyter.org}{Jupyter Notebook}.
\item
  \textbf{Google Colab:} para executar no seu navegador, vá em
  \textbf{Arquivo → Salvar uma cópia no Drive} e abra o arquivo
  \textbf{Cópia de aula01.ipynb}, que ficará na pasta \textbf{Colab
  Notebooks} do seu Drive.
\end{itemize}

O resultado será exibido logo abaixo da célula executada.

\begin{tcolorbox}[enhanced jigsaw, bottomtitle=1mm, bottomrule=.15mm, colback=white, colframe=quarto-callout-note-color-frame, rightrule=.15mm, opacityback=0, leftrule=.75mm, titlerule=0mm, arc=.35mm, toptitle=1mm, colbacktitle=quarto-callout-note-color!10!white, title=\textcolor{quarto-callout-note-color}{\faInfo}\hspace{0.5em}{Nota sobre o formato}, opacitybacktitle=0.6, left=2mm, toprule=.15mm, breakable, coltitle=black]

Você pode estar acessando este conteúdo de diferentes maneiras. Se
estiver executando este notebook no \href{https://jupyter.org/}{Jupyter}
ou no \href{https://colab.research.google.com/}{Google Colab}, as
células de código são interativas: você pode executá-las pressionando
Shift + Enter ou clicando no botão de \emph{play}. Se estiver lendo a
versão renderizada em HTML ou PDF, o código aparecerá como blocos
estáticos, mas o conteúdo e as explicações permanecem os mesmos --- a
única diferença é que você não poderá executar o código diretamente no
documento.

\end{tcolorbox}

\part{Parte I --- Fundamentos de PDI}

\chapter{Fundamentos e Primeiros
Passos}\label{fundamentos-e-primeiros-passos}

\href{https://colab.research.google.com/github/fzampirolli/pdi-vc/blob/master/notebooks_alunos/cap01/cap01_aluno.ipynb}{
\includegraphics[width=0.16\linewidth,height=\textheight,keepaspectratio]{cap01/imagens/colab-badge.png}}
\href{https://github.com/fzampirolli/pdi-vc}{
\includegraphics[width=0.2\linewidth,height=\textheight,keepaspectratio]{cap01/imagens/github-badge.png}}

Este capítulo inaugura a \textbf{Parte 1} do livro, focada nos
fundamentos do Processamento Digital de Imagens (PDI). Aqui, você
aprenderá como as imagens são representadas matematicamente e como
manipulá-las usando Python e a biblioteca \texttt{morph.py}
\autocite{zampirolli2025morph}.

\section{Objetivos}\label{objetivos}

Ao final deste capítulo, você será capaz de:

\begin{itemize}
\tightlist
\item
  Compreender a natureza física e matemática da imagem digital
  \(f(x,y)\).
\item
  Identificar as faixas do espectro eletromagnético relevantes para PDI.
\item
  Configurar o ambiente de desenvolvimento em Python.
\item
  Realizar operações básicas: leitura, exibição e salvamento de imagens.
\item
  Manipular estruturas de matrizes (NumPy) sem cair em armadilhas de
  memória.
\item
  Acessar e modificar intensidades de pixels individualmente.
\item
  Aplicar limiarização manual.
\end{itemize}

\section{\texorpdfstring{Antes de começar: \emph{Notebooks} em
Python}{Antes de começar: Notebooks em Python}}\label{antes-de-comeuxe7ar-notebooks-em-python-1}

Este material foi construído sob o conceito de \textbf{\emph{Literate
Programming}} (Programação Literária), idealizado por Donald Knuth na
década de 84 \autocite{knuth_literate_1984}. Knuth --- também criador do
sistema \textbf{TeX} para tipografia digital --- propôs que os programas
fossem escritos como uma narrativa lógica para seres humanos,
intercalando código e documentação.

Para executar uma célula, pressione Shift + Enter ou clique no botão ▶️.

\begin{tcolorbox}[enhanced jigsaw, bottomtitle=1mm, bottomrule=.15mm, colback=white, colframe=quarto-callout-note-color-frame, rightrule=.15mm, opacityback=0, leftrule=.75mm, titlerule=0mm, arc=.35mm, toptitle=1mm, colbacktitle=quarto-callout-note-color!10!white, title=\textcolor{quarto-callout-note-color}{\faInfo}\hspace{0.5em}{Nota sobre o formato}, opacitybacktitle=0.6, left=2mm, toprule=.15mm, breakable, coltitle=black]

Se você estiver lendo a versão renderizada (PDF ou HTML), o código
aparecerá como blocos estáticos. Para interagir, acesse o link do
\textbf{Google Colab} no topo da página.

\end{tcolorbox}

\section{Fundamentos}\label{fundamentos}

O estudo de sistemas baseados em imagens compreende um ecossistema de
disciplinas integradas que transformam dados visuais brutos em
conhecimento estruturado. Enquanto algumas áreas focam na geração de
representações, outras dedicam-se ao tratamento e à análise desses dados
para dar suporte a aplicações tecnológicas complexas.

O diagrama apresentado no Figura~\ref{fig-1-representacao} estabelece a
distinção e complementaridade entre o Processamento de Imagens (PDI) e a
Visão Computacional (VC). O PDI, destacado em \textbf{verde}, tem como
foco a transformação de imagem para imagem, visando melhoria de
qualidade ou pré-processamento, como a remoção de ruídos e realce de
contraste.

Em contrapartida, a VC, assinalada em \textbf{azul}, foca na
interpretação do conteúdo visual para extrair modelos ou informações,
como o reconhecimento de objetos e gestos. A região de intersecção
ilustra a sinergia entre as áreas, onde o PDI prepara os dados visuais
para a interpretação pela VC. O mapa também demonstra as interconexões
de ambas as disciplinas com áreas como Robótica, Computação Gráfica,
Inteligência Artificial e Neurociência.

\begin{figure}

\centering{

\includegraphics[width=1\linewidth,height=\textheight,keepaspectratio]{cap01/imagens/fig-1-pdf-vc.png}

}

\caption{\label{fig-1-representacao}Diagrama relacional detalhando as
distinções fundamentais, sinergias e interconexões entre Processamento
de Imagens e Visão Computacional no contexto de sistemas baseados em
imagens.}

\end{figure}%

\subsection{👁️ Visão Computacional}\label{visuxe3o-computacional}

\begin{itemize}
\tightlist
\item
  \textbf{Foco:} \emph{Imagem → Modelo} (caminho inverso da Computação
  Gráfica).\\
\item
  \textbf{Objetivo:} Extrair informação de alto nível a partir de
  imagens ou vídeos.\\
\item
  \textbf{Aplicações típicas:}

  \begin{itemize}
  \tightlist
  \item
    Robótica -- detecção de obstáculos, localização e navegação
    autônoma.
  \item
    Vigilância e inspeção -- reconhecimento de eventos, leitura de
    placas, controle de qualidade.
  \item
    Sensoriamento remoto -- análise de imagens de satélite, mapeamento
    ambiental.
  \item
    Imagens médicas -- detecção de tumores, segmentação de órgãos,
    auxílio ao diagnóstico.
  \item
    Interação humano‑computador -- reconhecimento de gestos, expressões
    faciais, rastreamento ocular.
  \end{itemize}
\item
  \textbf{Relação com outras áreas:} utiliza técnicas de Aprendizado de
  Máquina e IA para classificar e interpretar cenas; serve como
  ``olhos'' da Robótica.
\end{itemize}

\subsection{🖼️ Processamento Digital de Imagens
(PDI)}\label{processamento-digital-de-imagens-pdi}

\begin{itemize}
\tightlist
\item
  \textbf{Foco:} \emph{Imagem → Imagem} (geralmente -- transformação de
  uma imagem em outra).\\
\item
  \textbf{Objetivo:} Melhorar a qualidade visual ou extrair
  características de baixo nível.\\
\item
  \textbf{Aplicações comuns:}

  \begin{itemize}
  \tightlist
  \item
    Eliminação de ruídos (filtros de média, mediana, gaussiano).
  \item
    Melhoria de contraste (equalização de histograma, ajuste gamma).
  \item
    Detecção de bordas (Sobel, Canny, Laplaciano).
  \item
    Segmentação (limiarização, crescimento de regiões,
    \emph{watershed}).
  \item
    Transformações geométricas (redimensionamento, rotação, correção de
    perspectiva).
  \end{itemize}
\item
  \textbf{Relação com outras áreas:}

  \begin{itemize}
  \tightlist
  \item
    É a \textbf{base} para a maioria dos sistemas de Visão Computacional
    (pré‑processamento).\\
  \item
    A Computação Gráfica frequentemente aplica PDI para
    pós‑processamento (ex.: suavização, realce).\\
  \item
    Técnicas de IA podem otimizar parâmetros de processamento (ex.:
    aprendizado de filtros).
  \end{itemize}
\end{itemize}

\subsection{Como as ``Outras Áreas'' se
conectam}\label{como-as-outras-uxe1reas-se-conectam}

\begin{longtable}[]{@{}
  >{\raggedright\arraybackslash}p{(\linewidth - 2\tabcolsep) * \real{0.2389}}
  >{\raggedright\arraybackslash}p{(\linewidth - 2\tabcolsep) * \real{0.7611}}@{}}
\toprule\noalign{}
\begin{minipage}[b]{\linewidth}\raggedright
Área
\end{minipage} & \begin{minipage}[b]{\linewidth}\raggedright
Relação com PDI e VC
\end{minipage} \\
\midrule\noalign{}
\endhead
\bottomrule\noalign{}
\endlastfoot
\textbf{Inteligência Artificial} & Fornece modelos (redes neurais, SVM)
que interpretam saídas da VC. \\
\textbf{Robótica} & Consome dados de VC para tomar decisões (navegação,
manipulação). \\
\textbf{Aprendizado de Máquina} & Usa descritores extraídos pelo PDI/VC
para treinar classificadores. \\
\textbf{Computação Gráfica} & Caminho inverso: \emph{modelo → imagem};
muitas vezes aplica PDI para renderização realista. \\
\textbf{Neurociência} & Inspira modelos de PDI (ex.: filtros semelhantes
a células ganglionares da retina). \\
\end{longtable}

\section{Etapas do PDI}\label{etapas-do-pdi}

As etapas do PDI são apresentadas na Figura~\ref{fig-1-etapas}, que
podem ser compreendidas como uma cadeia de transformações que reduz a
redundância dos dados em busca de significado:

\begin{itemize}
\tightlist
\item
  \textbf{Baixo Nível:} Atua diretamente sobre os pixels da imagem
  ruidosa para realizar melhorias e filtragens, gerando como saída uma
  imagem limpa ou realçada.
\item
  \textbf{Médio Nível:} Recebe a imagem tratada e realiza a segmentação
  e descrição, transformando a matriz de pixels em atributos
  estruturados (forma, tamanho e textura).
\item
  \textbf{Alto Nível:} Utiliza a tabela de atributos para alimentar
  processos de lógica e inteligência artificial, resultando na decisão
  ou reconhecimento final (como o diagnóstico médico).
\end{itemize}

\begin{figure}

\centering{

\includegraphics[width=1\linewidth,height=\textheight,keepaspectratio]{cap01/imagens/fig-1-etapas.png}

}

\caption{\label{fig-1-etapas}\textbf{Representação do fluxo sequencial
de processamento: a saída de cada nível torna-se a entrada do nível
subsequente.}}

\end{figure}%

Detalhando mais essas etapas, a Figura~\ref{fig-1-etapas2} ilustra a
jornada completa da informação visual, desde a sua captura até a
interpretação final. O fluxo inicia-se na \textbf{Aquisição da Imagem
(1)} e avança por processos de tratamento como o \textbf{Melhoramento
(2)} e a \textbf{Restauração (3)}. Em seguida, o conteúdo é isolado
através da \textbf{Segmentação (4)} e refinado pela \textbf{Morfologia
(5)}. O estágio crucial de transição ocorre na \textbf{Representação e
Descrição (6)}, onde objetos visuais são convertidos em dados
matemáticos (como área e perímetro), permitindo que o sistema execute o
\textbf{Reconhecimento do Objeto (7)}. Este ecossistema é suportado por
pilares como o \textbf{Processamento de Imagem Colorida} e a
\textbf{Compressão}, que garantem a eficiência no armazenamento e
análise dos dados.

\begin{figure}

\centering{

\includegraphics[width=1\linewidth,height=\textheight,keepaspectratio]{cap01/imagens/fig-1-etapas2.png}

}

\caption{\label{fig-1-etapas2}\textbf{Fluxo detalhado do PDI: da
aquisição sensorial à extração de atributos e reconhecimento
automatizado, incluindo em processamento colorido e compressão.}}

\end{figure}%

\section{Formação da Imagem e o
Espectro}\label{formauxe7uxe3o-da-imagem-e-o-espectro}

O processo de formação de uma imagem é fundamentado na interação entre a
matéria e a energia radiante. Essencialmente, uma imagem é concebida
quando um \textbf{sensor registra a radiação} resultante da interação
com um objeto físico. No contexto da visão humana e da fotografia
convencional, este fenômeno depende de uma fonte de luz que ilumine a
cena; as características dos objetos são então codificadas através das
\textbf{variações de intensidade e cor} da luz que atinge o sensor,
conforme ilustrado na Figura~\ref{fig-1-luz}.

\begin{figure}

\centering{

\includegraphics[width=1\linewidth,height=\textheight,keepaspectratio]{cap01/imagens/fig-1-luz.png}

}

\caption{\label{fig-1-luz}Representação do espectro visível e sua
posição em relação às demais radiações eletromagnéticas, destacando a
variação dos comprimentos de onda de 380 nm a 750 nm.}

\end{figure}%

A \textbf{luz visível} ocupa apenas uma pequena faixa do espectro
eletromagnético --- entre \textbf{380 nm (violeta)} e \textbf{750 nm
(vermelho)} ---, conforme ilustrado na Figura~\ref{fig-1-espectro}.
Sensores digitais convencionais operam nessa mesma janela, mas
equipamentos especializados podem captar radiações invisíveis ao olho
humano, como o infravermelho e os raios X. Em PDI, a imagem formada
depende diretamente da sensibilidade espectral do sensor utilizado.

\begin{figure}

\centering{

\pandocbounded{\includegraphics[keepaspectratio]{cap01/cap01.py.pt_files/figure-pdf/fig-1-espectro-output-1.png}}

}

\caption{\label{fig-1-espectro}(A) Espectro eletromagnético completo em
escala logarítmica, com destaque para a faixa visível. (B) Detalhe da
luz visível (380--750 nm) e suas cores componentes. (C) Decomposição da
luz branca pelo prisma: menor comprimento de onda sofre maior refração,
separando UV, visível e infravermelho.}

\end{figure}%

\section{O que é uma Imagem
Digital?}\label{o-que-uxe9-uma-imagem-digital}

Uma \textbf{imagem digital} é formada por uma grade de \textbf{pixels}
(\emph{Picture Elements}), onde cada pixel é a menor unidade elementar
da imagem.

\begin{tcolorbox}[enhanced jigsaw, bottomtitle=1mm, bottomrule=.15mm, colback=white, colframe=quarto-callout-tip-color-frame, rightrule=.15mm, opacityback=0, leftrule=.75mm, titlerule=0mm, arc=.35mm, toptitle=1mm, colbacktitle=quarto-callout-tip-color!10!white, title=\textcolor{quarto-callout-tip-color}{\faLightbulb}\hspace{0.5em}{O que é um Pixel?}, opacitybacktitle=0.6, left=2mm, toprule=.15mm, breakable, coltitle=black]

Um \textbf{pixel} é a menor unidade que compõe uma imagem digital. Cada
pixel ocupa uma posição única na grade e armazena um ou mais valores
numéricos que representam sua intensidade ou cor.

\end{tcolorbox}

\subsection*{Representação
Matemática}\label{representauxe7uxe3o-matemuxe1tica}
\addcontentsline{toc}{subsection}{Representação Matemática}

Diferente de uma função contínua, o domínio de uma imagem digital é um
plano retangular finito \(\mathbb{E} \subset \mathbb{Z}^2\), que
representa a grade de amostragem. Este domínio é indexado por
coordenadas inteiras:

\begin{equation}\phantomsection\label{eq-1-dominio}{
\mathbb{E}  = \{ (x, y) \in \mathbb{Z}^2 \mid 0 \le x < L,\; 0 \le y < H \}
}\end{equation}

Onde:

\begin{itemize}
\tightlist
\item
  \textbf{\(L\)}: representa a largura da imagem (número de colunas).
\item
  \textbf{\(H\)}: representa a altura da imagem (número de linhas).
\end{itemize}

A imagem digital é uma função que associa cada par de coordenadas
\((x,y)\) a um ou mais valores que descrevem a aparência do pixel.

\begin{equation}\phantomsection\label{eq-1-image}{
f: \mathbb{E}  \to \mathcal{V}
}\end{equation}

O conjunto \(\mathcal{V}\) define os valores possíveis para o pixel
(codomínio), variando conforme o tipo de imagem, como demonstrado na
Tabela~\ref{tbl-1-tipos-imagem}.

\begin{longtable}[]{@{}llc@{}}
\caption{Principais tipos de imagem digital e seus respectivos conjuntos
de valores possíveis para cada
pixel.}\label{tbl-1-tipos-imagem}\tabularnewline
\toprule\noalign{}
Tipo de imagem & \(\mathcal{V}\) (valores do pixel) & Representação \\
\midrule\noalign{}
\endfirsthead
\toprule\noalign{}
Tipo de imagem & \(\mathcal{V}\) (valores do pixel) & Representação \\
\midrule\noalign{}
\endhead
\bottomrule\noalign{}
\endlastfoot
\textbf{Binária} & \(\{0, 1\}\) ou \(\{0, 255\}\) & ⬛◻️ \\
\textbf{Tons de cinza} & \(\{0, 1, \dots, 255\}\) & ░▒▓█ \\
\textbf{Colorida (RGB)} & \(\{0, \dots, 255\}^3\) (triplas) & 🟥🟩🟦 \\
\end{longtable}

\textbf{Exemplo prático:} Uma imagem colorida (RGB) é representada
matematicamente por uma função que retorna três valores para cada pixel.
No computador, isso resulta em três matrizes sobrepostas (canais R, G e
B), onde cada célula contém um valor de intensidade para aquele canal
específico no pixel.

\section{Configuração do Ambiente}\label{configurauxe7uxe3o-do-ambiente}

Utilizaremos o ecossistema científico do Python, com destaque para o
\textbf{NumPy} (matemática matricial), \textbf{OpenCV} (visão
computacional padrão de mercado) e a biblioteca \textbf{morph.py}
\autocite{zampirolli2025morph}, desenvolvida especificamente para fins
didáticos neste livro.

\subsection{📦 Bibliotecas utilizadas}\label{bibliotecas-utilizadas}

\begin{longtable}[]{@{}
  >{\raggedright\arraybackslash}p{(\linewidth - 2\tabcolsep) * \real{0.4000}}
  >{\raggedright\arraybackslash}p{(\linewidth - 2\tabcolsep) * \real{0.6000}}@{}}
\toprule\noalign{}
\begin{minipage}[b]{\linewidth}\raggedright
Biblioteca
\end{minipage} & \begin{minipage}[b]{\linewidth}\raggedright
Função principal
\end{minipage} \\
\midrule\noalign{}
\endhead
\bottomrule\noalign{}
\endlastfoot
\texttt{numpy} & Representação matricial de imagens \\
\texttt{opencv-python} & Leitura, escrita e operações de visão
computacional \\
\texttt{matplotlib} & Visualização de imagens e gráficos \\
\texttt{morph.py} & Abstração didática das operações de PDI \\
\end{longtable}

\begin{Shaded}
\begin{Highlighting}[]
\CommentTok{\# Instalação (caso necessário no ambiente local)}
\CommentTok{\# !pip install opencv{-}python numpy matplotlib scikit{-}image gdown}

\ImportTok{import}\NormalTok{ sys}
\ImportTok{import}\NormalTok{ importlib}
\ImportTok{import}\NormalTok{ numpy }\ImportTok{as}\NormalTok{ np}
\ImportTok{import}\NormalTok{ cv2}
\ImportTok{import}\NormalTok{ matplotlib.pyplot }\ImportTok{as}\NormalTok{ plt}

\CommentTok{\# Verificar onde esté morph.py para importação}
\NormalTok{sys.path.insert(}\DecValTok{0}\NormalTok{, }\StringTok{\textquotesingle{}../../morph\textquotesingle{}}\NormalTok{)}
\ImportTok{import}\NormalTok{ morph}
\NormalTok{importlib.}\BuiltInTok{reload}\NormalTok{(morph)}
\ImportTok{from}\NormalTok{ morph }\ImportTok{import}\NormalTok{ mm}

\BuiltInTok{print}\NormalTok{(}\SpecialStringTok{f"Ambiente configurado:"}\NormalTok{)}
\BuiltInTok{print}\NormalTok{(}\SpecialStringTok{f"{-} OpenCV: }\SpecialCharTok{\{}\NormalTok{cv2}\SpecialCharTok{.}\NormalTok{\_\_version\_\_}\SpecialCharTok{\}}\SpecialStringTok{"}\NormalTok{)}
\BuiltInTok{print}\NormalTok{(}\SpecialStringTok{f"{-} NumPy: }\SpecialCharTok{\{}\NormalTok{np}\SpecialCharTok{.}\NormalTok{\_\_version\_\_}\SpecialCharTok{\}}\SpecialStringTok{"}\NormalTok{)}
\BuiltInTok{help}\NormalTok{(mm.show)}
\end{Highlighting}
\end{Shaded}

\begin{verbatim}
Ambiente configurado:
- OpenCV: 4.12.0
- NumPy: 2.2.6
Help on function show in module morph:

show(*args, title=None)
    This function will draw images f
    input: <*args> set of images f_i, where i>0 is binary image
          <title> optional string title for the plot
    output: image drawing
    Example:
    f1, f2 = np.zeros((100, 100,3)),  np.zeros((100, 100))
    f2[50:60, 50:60] = 1
    mm.show(f1, f2, title='Exemplo')
\end{verbatim}

\section{Fundamentos de Matrizes
(Cuidado!)}\label{fundamentos-de-matrizes-cuidado}

Como uma imagem digital é uma matriz, precisamos saber criá-las
corretamente. Em Python, existe uma armadilha comum ao usar o operador *
em listas:

\begin{tcolorbox}[enhanced jigsaw, bottomtitle=1mm, bottomrule=.15mm, colback=white, colframe=quarto-callout-warning-color-frame, rightrule=.15mm, opacityback=0, leftrule=.75mm, titlerule=0mm, arc=.35mm, toptitle=1mm, colbacktitle=quarto-callout-warning-color!10!white, title=\textcolor{quarto-callout-warning-color}{\faExclamationTriangle}\hspace{0.5em}{Atenção à cópia de referências}, opacitybacktitle=0.6, left=2mm, toprule=.15mm, breakable, coltitle=black]

Ao fazer \texttt{m\ =\ {[}{[}0{]}*2{]}*3}, você não cria 3 linhas
independentes, mas sim 3 referências para a \textbf{mesma} linha.
Alterar um valor em uma linha alterará todas as outras!

\end{tcolorbox}

Para visualizar esse comportamento, você pode testar o código no
\href{https://pythontutor.com/}{Python Tutor} e comparar com a forma
correta de criar matriz com listas:
\texttt{m\ =\ {[}{[}0{]}*2\ for\ \_\ in\ range(3){]}}.

A forma recomendada e padrão em Processamento Digital de Imagens é usar
\textbf{NumPy}, que cria matrizes com dados independentes e oferece
eficiência computacional. O código a seguir mostra diferentes formas de
criação de imagens sintéticas --- o resultado é exibido na
Figura~\ref{fig-imagens-sinteticas}.

\begin{Shaded}
\begin{Highlighting}[]
\CommentTok{\# Criando uma imagem preta (zeros) de 5x5 pixels}
\NormalTok{img\_preta }\OperatorTok{=}\NormalTok{ np.zeros((}\DecValTok{5}\NormalTok{, }\DecValTok{5}\NormalTok{), dtype}\OperatorTok{=}\StringTok{\textquotesingle{}uint8\textquotesingle{}}\NormalTok{)}

\CommentTok{\# Criando uma imagem branca (255) de 5x5 pixels}
\NormalTok{img\_branca }\OperatorTok{=}\NormalTok{ np.ones((}\DecValTok{5}\NormalTok{, }\DecValTok{5}\NormalTok{), dtype}\OperatorTok{=}\StringTok{\textquotesingle{}uint8\textquotesingle{}}\NormalTok{) }\OperatorTok{*} \DecValTok{255}

\CommentTok{\# Criando uma imagem aleatória para testes (Ruído)}
\NormalTok{img\_random }\OperatorTok{=}\NormalTok{ mm.randomImage(}\DecValTok{4}\NormalTok{, }\DecValTok{6}\NormalTok{, maxValue}\OperatorTok{=}\DecValTok{255}\NormalTok{)}

\BuiltInTok{print}\NormalTok{(}\StringTok{"Matriz Aleatória Gerada:"}\NormalTok{)}
\BuiltInTok{print}\NormalTok{(mm.drawImage(img\_random))}

\NormalTok{mm.show(img\_random, title}\OperatorTok{=}\StringTok{"Exemplo: Captura RGB"}\NormalTok{)}
\end{Highlighting}
\end{Shaded}

\begin{verbatim}
Matriz Aleatória Gerada:
161  90 244 168 122 169 
208  44  93 105 120  59 
 84  82 144 172  48 249 
121  94  76 251 170 107 
\end{verbatim}

\begin{figure}[H]

\centering{

\pandocbounded{\includegraphics[keepaspectratio]{cap01/cap01.py.pt_files/figure-pdf/fig-imagens-sinteticas-output-2.png}}

}

\caption{\label{fig-imagens-sinteticas}Exemplo de imagem sintética
gerada por matriz aleatória 4x6.}

\end{figure}%

\section{Lendo e Exibindo Imagens}\label{lendo-e-exibindo-imagens}

Em Python, as imagens são lidas como \emph{arrays} do NumPy. Isso
significa que toda a potência da álgebra linear está disponível para o
processamento de imagens.

A operação mais básica em PDI é a leitura de uma imagem. A função
\texttt{mm.read()} da \texttt{morph.py} aceita caminhos locais e URLs,
ver o resultado na Figura~\ref{fig-1-lena}.

\begin{Shaded}
\begin{Highlighting}[]
\CommentTok{\# Carregando uma imagem de exemplo via URL}
\NormalTok{url }\OperatorTok{=} \StringTok{"https://upload.wikimedia.org/wikipedia/en/7/7d/Lenna\_\%28test\_image\%29.png"}
\NormalTok{img }\OperatorTok{=}\NormalTok{ mm.read(url)}

\CommentTok{\# Exibindo informações básicas}
\BuiltInTok{print}\NormalTok{(}\SpecialStringTok{f"Dimensões (H, W, Canais): }\SpecialCharTok{\{}\NormalTok{img}\SpecialCharTok{.}\NormalTok{shape}\SpecialCharTok{\}}\SpecialStringTok{"}\NormalTok{)}
\BuiltInTok{print}\NormalTok{(}\SpecialStringTok{f"Tipo de dado: }\SpecialCharTok{\{}\NormalTok{img}\SpecialCharTok{.}\NormalTok{dtype}\SpecialCharTok{\}}\SpecialStringTok{"}\NormalTok{)}

\CommentTok{\# Exibição estética}
\NormalTok{mm.show(img, title}\OperatorTok{=}\StringTok{"Exemplo: Captura RGB"}\NormalTok{)}
\end{Highlighting}
\end{Shaded}

\begin{verbatim}
Dimensões (H, W, Canais): (512, 512, 3)
Tipo de dado: uint8
\end{verbatim}

\begin{figure}[H]

\centering{

\pandocbounded{\includegraphics[keepaspectratio]{cap01/cap01.py.pt_files/figure-pdf/fig-1-lena-output-2.png}}

}

\caption{\label{fig-1-lena}Imagem clássica de processamento de imagens:
`Lena'.}

\end{figure}%

\subsection*{Alternativa: Baixando a imagem
localmente}\label{alternativa-baixando-a-imagem-localmente}
\addcontentsline{toc}{subsection}{Alternativa: Baixando a imagem
localmente}

Se por algum motivo a leitura direta da URL não funcionar (firewall,
restrições de rede, ou se você preferir trabalhar com arquivos locais),
uma abordagem simples é \textbf{baixar a imagem} usando o comando
\texttt{wget} (disponível em ambientes Linux, macOS e no Google Colab) e
depois carregá-la com \texttt{mm.read()} a partir do arquivo salvo.

\begin{Shaded}
\begin{Highlighting}[]
\ExtensionTok{!wget} \AttributeTok{{-}O}\NormalTok{ lena.png https://upload.wikimedia.org/wikipedia/en/7/7d/Lenna\_\%28test\_image\%29.png}
\end{Highlighting}
\end{Shaded}

\begin{Shaded}
\begin{Highlighting}[]
\NormalTok{img }\OperatorTok{=}\NormalTok{ mm.read(}\StringTok{\textquotesingle{}lena.png\textquotesingle{}}\NormalTok{)}
\NormalTok{mm.show(img, title}\OperatorTok{=}\StringTok{"Exemplo: Captura RGB (arquivo local)"}\NormalTok{)}
\end{Highlighting}
\end{Shaded}

\subsubsection*{Explicação:}\label{explicauxe7uxe3o}
\addcontentsline{toc}{subsubsection}{Explicação:}

\begin{itemize}
\tightlist
\item
  \texttt{!wget\ -O\ lena.png\ \textless{}URL\textgreater{}} baixa a
  imagem da Wikimedia e a salva com o nome \texttt{lena.png}.
\item
  \texttt{mm.read(\textquotesingle{}lena.png\textquotesingle{})} lê o
  arquivo local (sem necessidade de cabeçalhos HTTP ou tratamento
  especial).
\item
  Essa estratégia é útil quando você quer \textbf{evitar dependências de
  rede} durante a execução ou \textbf{reutilizar a mesma imagem} várias
  vezes sem novo download.
\end{itemize}

\begin{tcolorbox}[enhanced jigsaw, bottomtitle=1mm, bottomrule=.15mm, colback=white, colframe=quarto-callout-tip-color-frame, rightrule=.15mm, opacityback=0, leftrule=.75mm, titlerule=0mm, arc=.35mm, toptitle=1mm, colbacktitle=quarto-callout-tip-color!10!white, title=\textcolor{quarto-callout-tip-color}{\faLightbulb}\hspace{0.5em}{Nota}, opacitybacktitle=0.6, left=2mm, toprule=.15mm, breakable, coltitle=black]

O prefixo \texttt{!} funciona no Jupyter Notebook, JupyterLab e Google
Colab. Em um terminal comum, basta remover o \texttt{!}. Se o
\texttt{wget} não estiver instalado, instale-o com
\texttt{apt\ install\ wget} (Linux) ou use
\texttt{curl\ -o\ lena.png\ \textless{}URL\textgreater{}} como
alternativa.

\end{tcolorbox}

\section{Conversão de Tipos e
Limiarização}\label{conversuxe3o-de-tipos-e-limiarizauxe7uxe3o}

Conforme vimos, uma imagem colorida no espaço RGB é representada pela
função:

\[
f: \mathbb{E} \to \{0,1,\dots,255\}^3
\]

Ou seja, para cada pixel \((x,y)\), temos três valores \((R,G,B)\) que
definem sua cor.

\subsection*{Conversão para Tons de
Cinza}\label{conversuxe3o-para-tons-de-cinza}
\addcontentsline{toc}{subsection}{Conversão para Tons de Cinza}

Para converter uma imagem RGB em tons de cinza (\emph{grayscale}), é
necessário combinar os três canais em um único valor de intensidade
\(g\), que representa o brilho percebido. Como o olho humano não é
igualmente sensível ao vermelho, verde e azul, utiliza-se uma
\textbf{média ponderada}. O padrão
\href{https://www.itu.int/rec/R-REC-BT.601/}{ITU-R BT.601}
\autocite{itu_bt601} define os seguintes pesos:

\begin{equation}\phantomsection\label{eq-1-cinza}{
g = 0.299\,R + 0.587\,G + 0.114\,B 
}\end{equation}

Após o cálculo, o valor \(g\) é arredondado para o inteiro mais próximo
e ajustado ao intervalo \([0, 255]\). O resultado é uma nova imagem,
agora em tons de cinza, representada por:

\[
f_{\text{cinza}}: \mathbb{E} \to \{0,1,\dots,255\}
\]

\subsection*{\texorpdfstring{Limiarização
(\emph{Thresholding})}{Limiarização (Thresholding)}}\label{limiarizauxe7uxe3o-thresholding}
\addcontentsline{toc}{subsection}{Limiarização (\emph{Thresholding})}

A partir da imagem em tons de cinza \(f_{\text{cinza}}(x,y)\), uma
operação fundamental é a \textbf{limiarização}, que produz uma imagem
\textbf{binária} (apenas preto e branco). Para isso, escolhe-se um valor
de corte \(T\) (geralmente no intervalo \([0,255]\)) e define-se:

\begin{equation}\phantomsection\label{eq-1-limiar}{
f_{\text{bin}}(x,y) = 
\begin{cases} 
255 & \text{se } f_{\text{cinza}}(x,y) > T \\[4pt]
0 & \text{caso contrário}
\end{cases}
}\end{equation}

\textbf{Exemplo:} Com \(T = 128\), pixels com intensidade acima de 128
tornam-se brancos (255); os demais tornam-se pretos (0).

A limiarização é amplamente usada para \textbf{segmentar objetos do
fundo}, extrair bordas ou criar máscaras binárias para processamento
posterior.

\begin{quote}
\textbf{Nota:} O valor 255 representa o branco máximo em imagens de 8
bits, enquanto 0 representa o preto absoluto.
\end{quote}

\subsection*{Exemplo prático de conversão e
limiarização}\label{exemplo-pruxe1tico-de-conversuxe3o-e-limiarizauxe7uxe3o}
\addcontentsline{toc}{subsection}{Exemplo prático de conversão e
limiarização}

A Figura~\ref{fig-1-processamento-basico} ilustra os principais passos
para transformar uma imagem colorida em tons de cinza e, em seguida,
binarizá‑la por limiarização. O código a seguir implementa essas etapas:

\begin{Shaded}
\begin{Highlighting}[]
\ImportTok{import}\NormalTok{ matplotlib.pyplot }\ImportTok{as}\NormalTok{ plt}

\CommentTok{\# 1. Converter para Tons de Cinza}
\NormalTok{img\_gray }\OperatorTok{=}\NormalTok{ mm.gray(img)}

\CommentTok{\# 2. Aplicar limiar (Pixels \textgreater{} 128 tornam{-}se 255, outros 0)}
\NormalTok{limiar }\OperatorTok{=} \DecValTok{128}
\NormalTok{img\_binaria }\OperatorTok{=}\NormalTok{ mm.threshold(img\_gray, limiar)}

\CommentTok{\# Uso da nova função}
\NormalTok{mm.show\_multiple(}
\NormalTok{    [img, img\_gray, img\_binaria], }
\NormalTok{    titles}\OperatorTok{=}\NormalTok{[}\StringTok{"Original"}\NormalTok{, }\StringTok{"Tons de Cinza"}\NormalTok{, }\SpecialStringTok{f"Binária (T=}\SpecialCharTok{\{}\NormalTok{limiar}\SpecialCharTok{\}}\SpecialStringTok{)"}\NormalTok{],}
\NormalTok{    cols}\OperatorTok{=}\DecValTok{3}
\NormalTok{)}
\end{Highlighting}
\end{Shaded}

\begin{figure}[H]

\centering{

\pandocbounded{\includegraphics[keepaspectratio]{cap01/cap01.py.pt_files/figure-pdf/fig-1-processamento-basico-output-1.png}}

}

\caption{\label{fig-1-processamento-basico}Processamento básico de
imagens: (a) imagem original, (b) imagem em tons de cinza, (c) imagem
binarizada por limiar (T=128).}

\end{figure}%

\section{Limiarização pelo método de
Otsu}\label{limiarizauxe7uxe3o-pelo-muxe9todo-de-otsu}

Conforme apresentado na Equação~\ref{eq-1-limiar}, a limiarização
converte uma imagem em tons de cinza para binária usando um valor de
corte \(T\). Até agora fixamos \(T = 128\) manualmente.

\begin{Shaded}
\begin{Highlighting}[]
\NormalTok{img\_bin\_fixo }\OperatorTok{=}\NormalTok{ mm.threshold(img\_gray, T}\OperatorTok{=}\DecValTok{128}\NormalTok{)}
\end{Highlighting}
\end{Shaded}

Entretanto, a escolha manual de \(T\) nem sempre é trivial. A biblioteca
\texttt{mm} oferece uma alternativa automática: quando o parâmetro
\texttt{limiar} \textbf{não é fornecido}, a função
\texttt{mm.threshold(img\_gray)} calcula o valor de \(T\) pelo
\textbf{método de Otsu} \autocite{otsu1979}. Esse método, que será
detalhado em capítulos futuros, maximiza a variância inter‑classes do
histograma, separando automaticamente os pixels de objeto e fundo.

O código abaixo compara a limiarização manual (\(T=128\)) com a
automática (Otsu), mostrando também o valor de \(T\) calculado:

\begin{Shaded}
\begin{Highlighting}[]
\CommentTok{\# Limiarização com T fixo (manual)}
\NormalTok{T\_fixo }\OperatorTok{=} \DecValTok{128}
\NormalTok{img\_bin\_fixo }\OperatorTok{=}\NormalTok{ mm.threshold(img\_gray, T\_fixo)}

\CommentTok{\# Limiarização pelo método de Otsu (T automático)}
\NormalTok{T\_otsu, img\_bin\_otsu }\OperatorTok{=}\NormalTok{ cv2.threshold(img\_gray, }\DecValTok{0}\NormalTok{, }\DecValTok{255}\NormalTok{,}
\NormalTok{                                      cv2.THRESH\_BINARY }\OperatorTok{+}\NormalTok{ cv2.THRESH\_OTSU)}
\BuiltInTok{print}\NormalTok{(}\SpecialStringTok{f\textquotesingle{}Limiar calculado por Otsu: T = }\SpecialCharTok{\{}\NormalTok{T\_otsu}\SpecialCharTok{\}}\SpecialStringTok{\textquotesingle{}}\NormalTok{)}
\CommentTok{\# ou simplesmente:}
\CommentTok{\# img\_bin\_otsu = mm.threshold(img\_gray)}

\CommentTok{\# Exibição lado a lado}
\NormalTok{mm.show\_multiple(}
\NormalTok{    [img\_gray, img\_bin\_fixo, img\_bin\_otsu], }
\NormalTok{    titles}\OperatorTok{=}\NormalTok{[}\StringTok{"Tons de Cinza"}\NormalTok{, }\SpecialStringTok{f"Binária (T=}\SpecialCharTok{\{}\NormalTok{T\_fixo}\SpecialCharTok{\}}\SpecialStringTok{)"}\NormalTok{, }\SpecialStringTok{f"Binária (Otsu, T=}\SpecialCharTok{\{}\NormalTok{T\_otsu}\SpecialCharTok{\}}\SpecialStringTok{)"}\NormalTok{],}
\NormalTok{    cols}\OperatorTok{=}\DecValTok{3}
\NormalTok{)}
\end{Highlighting}
\end{Shaded}

\begin{verbatim}
Limiar calculado por Otsu: T = 117.0
\end{verbatim}

\begin{figure}[H]

\centering{

\pandocbounded{\includegraphics[keepaspectratio]{cap01/cap01.py.pt_files/figure-pdf/fig-1-otsu-output-2.png}}

}

\caption{\label{fig-1-otsu}Comparação entre limiarização manual (T=128)
e automática (Otsu) sobre a imagem em tons de cinza.}

\end{figure}%

A Figura~\ref{fig-1-otsu} mostra que o limiar obtido por Otsu se adapta
automaticamente à imagem, resultando em uma binarização mais eficiente
do que um valor fixo, especialmente quando as intensidades do objeto e
do fundo são bem separadas no histograma. Essa técnica é amplamente
utilizada em sistemas de visão computacional para binarização de
documentos, detecção de objetos e pré‑processamento de imagens.

\begin{tcolorbox}[enhanced jigsaw, bottomtitle=1mm, bottomrule=.15mm, colback=white, colframe=quarto-callout-tip-color-frame, rightrule=.15mm, opacityback=0, leftrule=.75mm, titlerule=0mm, arc=.35mm, toptitle=1mm, colbacktitle=quarto-callout-tip-color!10!white, title=\textcolor{quarto-callout-tip-color}{\faLightbulb}\hspace{0.5em}{Simplicidade da biblioteca \texttt{morph.py}}, opacitybacktitle=0.6, left=2mm, toprule=.15mm, breakable, coltitle=black]

Enquanto o OpenCV exige a chamada completa:

\begin{Shaded}
\begin{Highlighting}[]
\NormalTok{T\_otsu, img\_bin }\OperatorTok{=}\NormalTok{ cv2.threshold(img\_gray, }\DecValTok{0}\NormalTok{, }\DecValTok{255}\NormalTok{,}
\NormalTok{                                cv2.THRESH\_BINARY }\OperatorTok{+}\NormalTok{ cv2.THRESH\_OTSU)}
\end{Highlighting}
\end{Shaded}

a biblioteca \texttt{mm} abstrai toda essa complexidade: basta chamar
\texttt{mm.threshold(img\_gray)}. O limiar de Otsu é calculado
automaticamente e a imagem binária é retornada diretamente. Essa
abordagem permite que você se concentre no conceito, não nos detalhes de
implementação.

\end{tcolorbox}

\section{Acesso a Pixels}\label{acesso-a-pixels}

Em Python com NumPy, uma imagem é representada como um \textbf{array
multidimensional}. Para acessar um pixel específico, utilizam-se as
coordenadas \textbf{linha} (eixo Y) e \textbf{coluna} (eixo X):
\texttt{img{[}linha,\ coluna{]}}.

O código abaixo demonstra como obter o valor de um pixel em uma imagem
RGB e em sua versão em tons de cinza, além de criar uma pequena imagem
sintética para visualizar a estrutura de uma matriz de pixels.

\begin{Shaded}
\begin{Highlighting}[]
\ImportTok{import}\NormalTok{ numpy }\ImportTok{as}\NormalTok{ np}
\ImportTok{import}\NormalTok{ matplotlib.pyplot }\ImportTok{as}\NormalTok{ plt}

\CommentTok{\# Coordenadas do pixel que queremos examinar}
\NormalTok{r, c }\OperatorTok{=} \DecValTok{100}\NormalTok{, }\DecValTok{100}

\CommentTok{\# Acessa o pixel na imagem em tons de cinza (um valor escalar)}
\NormalTok{pixel\_cinza }\OperatorTok{=}\NormalTok{ img\_gray[r, c]}

\CommentTok{\# Acessa o pixel na imagem colorida RGB (vetor de 3 valores)}
\NormalTok{pixel\_rgb }\OperatorTok{=}\NormalTok{ img[r, c]}

\BuiltInTok{print}\NormalTok{(}\SpecialStringTok{f\textquotesingle{}Pixel na posição (}\SpecialCharTok{\{}\NormalTok{r}\SpecialCharTok{\}}\SpecialStringTok{,}\SpecialCharTok{\{}\NormalTok{c}\SpecialCharTok{\}}\SpecialStringTok{):\textquotesingle{}}\NormalTok{)}
\BuiltInTok{print}\NormalTok{(}\SpecialStringTok{f\textquotesingle{}  {-} Tons de cinza : }\SpecialCharTok{\{}\NormalTok{pixel\_cinza}\SpecialCharTok{\}}\SpecialStringTok{\textquotesingle{}}\NormalTok{)}
\BuiltInTok{print}\NormalTok{(}\SpecialStringTok{f\textquotesingle{}  {-} RGB           : R=}\SpecialCharTok{\{}\NormalTok{pixel\_rgb[}\DecValTok{0}\NormalTok{]}\SpecialCharTok{\}}\SpecialStringTok{, G=}\SpecialCharTok{\{}\NormalTok{pixel\_rgb[}\DecValTok{1}\NormalTok{]}\SpecialCharTok{\}}\SpecialStringTok{, B=}\SpecialCharTok{\{}\NormalTok{pixel\_rgb[}\DecValTok{2}\NormalTok{]}\SpecialCharTok{\}}\SpecialStringTok{\textquotesingle{}}\NormalTok{)}

\CommentTok{\# Cria uma imagem sintética 5×5 com todos os pixels pretos (0)}
\NormalTok{syn }\OperatorTok{=}\NormalTok{ np.zeros((}\DecValTok{5}\NormalTok{, }\DecValTok{5}\NormalTok{), dtype}\OperatorTok{=}\NormalTok{np.uint8)}
\CommentTok{\# Torna o pixel central (linha 2, coluna 2) branco (255)}
\NormalTok{syn[}\DecValTok{2}\NormalTok{, }\DecValTok{2}\NormalTok{] }\OperatorTok{=} \DecValTok{255}

\BuiltInTok{print}\NormalTok{(}\StringTok{"}\CharTok{\textbackslash{}n}\StringTok{Matriz da imagem sintética 5×5:"}\NormalTok{)}
\BuiltInTok{print}\NormalTok{(syn)}

\CommentTok{\# Exibe a imagem sintética}
\CommentTok{\# plt.figure(figsize=(3, 3))}
\CommentTok{\# plt.imshow(syn, cmap=\textquotesingle{}gray\textquotesingle{}, vmin=0, vmax=255)}
\CommentTok{\# plt.title("Imagem sintética 5×5\textbackslash{}n(pixel central branco)")}
\CommentTok{\# plt.axis(\textquotesingle{}off\textquotesingle{})}
\CommentTok{\# plt.tight\_layout()}
\CommentTok{\# plt.show()}
\CommentTok{\# ou simplesmente:}

\NormalTok{mm.show(syn, title}\OperatorTok{=}\StringTok{"Exemplo: Imagem Sintética 5x5"}\NormalTok{)}
\end{Highlighting}
\end{Shaded}

\begin{verbatim}
Pixel na posição (100,100):
  - Tons de cinza : 108
  - RGB           : R=182, G=74, B=87

Matriz da imagem sintética 5×5:
[[  0   0   0   0   0]
 [  0   0   0   0   0]
 [  0   0 255   0   0]
 [  0   0   0   0   0]
 [  0   0   0   0   0]]
\end{verbatim}

\begin{figure}[H]

\centering{

\pandocbounded{\includegraphics[keepaspectratio]{cap01/cap01.py.pt_files/figure-pdf/fig-1-pixel-acesso-output-2.png}}

}

\caption{\label{fig-1-pixel-acesso}Exemplo de acesso a pixels: (a)
imagem original com pixel destacado, (b) imagem em tons de cinza, (c)
matriz sintética 5×5 com pixel central branco.}

\end{figure}%

\textbf{Explicação linha a linha:}

\begin{enumerate}
\def\labelenumi{\arabic{enumi}.}
\tightlist
\item
  \texttt{img\_gray{[}r,\ c{]}} -- retorna um único número inteiro entre
  0 e 255, correspondente ao nível de cinza naquela posição.\\
\item
  \texttt{img{[}r,\ c{]}} -- retorna uma tupla ou array com três valores
  (R, G, B).\\
\item
  \texttt{np.zeros((5,5),\ dtype=np.uint8)} -- cria uma matriz
  \(5\times 5\) preenchida com zeros (preto).\\
\item
  \texttt{syn{[}2,\ 2{]}\ =\ 255} -- altera o elemento central da matriz
  para 255 (branco), demonstrando como modificar um pixel.\\
\item
  A exibição com \texttt{plt.imshow} revela um pequeno quadrado com um
  ponto branco no meio, ilustrando visualmente a estrutura discreta da
  imagem.
\end{enumerate}

A Figura~\ref{fig-1-pixel-acesso} mostra os valores impressos e a imagem
sintética. Perceba que a indexação em Python é \textbf{zero‑based}: o
canto superior esquerdo é \texttt{(0,0)}. A ordem \textbf{linha ×
coluna} corresponde à estrutura matricial da imagem: primeira dimensão →
altura (Y), segunda dimensão → largura (X).

\section{Resumo}\label{resumo}

Neste capítulo foram apresentados os fundamentos de representação de
imagens digitais:

\begin{itemize}
\tightlist
\item
  \textbf{Imagem digital} = função \(f(x,y)\) que mapeia coordenadas
  para intensidades (escalares ou vetoriais).
\item
  \textbf{Domínio}: conjunto finito
  \(\mathbb{E} = \{(x,y) \in \mathbb{Z}^2 \mid 0 \le x < L,\; 0 \le y < H\}\).
\item
  \textbf{Tipos principais}: binária (\(\mathcal{V} = \{0, 255\}\)),
  tons de cinza (\(\mathcal{V} = [0,255]\)) e RGB
  (\(\mathcal{V} = [0,255]^3\)).
\item
  A biblioteca \texttt{morph.py} (ou \texttt{mm}) oferece funções
  didáticas para operações básicas de PDI, como \texttt{mm.gray()},
  \texttt{mm.threshold()}, \texttt{mm.show\_multiple()}.
\item
  \textbf{Armadilha NumPy}: nunca use \texttt{{[}{[}0{]}*n{]}*m} para
  criar matrizes --- use sempre \texttt{np.zeros()} ou
  \texttt{np.ones()}.
\item
  \textbf{Limiarização} converte tons de cinza em binária; o
  \textbf{método de Otsu} determina o valor de corte automaticamente
  maximizando a variância inter-classes.
\item
  Acesso a pixels via \texttt{img{[}linha,\ coluna{]}}, com indexação
  \textbf{zero-based}.
\end{itemize}

O Capítulo 2 abordará \textbf{histogramas e equalização de contraste}.

\section{🤖 Uso do NotebookLM como Tutor
Complementar}\label{uso-do-notebooklm-como-tutor-complementar}

Nesta edição, além dos notebooks interativos no Google Colab,
incentivamos o uso do \textbf{NotebookLM} como ferramenta complementar
de aprendizagem. Essa ferramenta de IA utiliza exclusivamente os
documentos fornecidos pelo autor como base de conhecimento, garantindo
respostas alinhadas ao conteúdo do livro.

Para cada capítulo, preparamos um projeto específico na plataforma. Para
uma experiência de estudo ampliada, utilize o acesso abaixo:

\begin{tcolorbox}[enhanced jigsaw, bottomtitle=1mm, bottomrule=.15mm, colback=white, colframe=quarto-callout-important-color-frame, rightrule=.15mm, opacityback=0, leftrule=.75mm, titlerule=0mm, arc=.35mm, toptitle=1mm, colbacktitle=quarto-callout-important-color!10!white, title={🎓 Estude com o Tutor Inteligente}, opacitybacktitle=0.6, left=2mm, toprule=.15mm, breakable, coltitle=black]

Para interagir com o conteúdo deste capítulo, acesse o link a seguir. O
ambiente contém materiais didáticos em diferentes formatos, gerados a
partir do \textbf{PDF} do capítulo. Na plataforma, explore especialmente
as opções \textbf{Guia de Estudo} e \textbf{Conversa} para aprofundar
sua compreensão.

\href{https://notebooklm.google.com/notebook/aca1138a-02aa-4f98-b777-1e25795ca635}{🚀
ACESSAR NOTEBOOKLM: CAPÍTULO 01}

\end{tcolorbox}

\subsection*{Funcionalidades Disponíveis na
Plataforma}\label{funcionalidades-disponuxedveis-na-plataforma}
\addcontentsline{toc}{subsection}{Funcionalidades Disponíveis na
Plataforma}

O \textbf{NotebookLM} oferece uma suíte avançada de ferramentas baseadas
em IA para transformar o conteúdo estático do livro em uma experiência
de aprendizado dinâmica e multimídia. Conforme ilustrado na
Figura~\ref{fig-notebooklm}, a plataforma utiliza técnicas de
\textbf{RAG} (\emph{Retrieval-Augmented Generation}), fundamentadas no
trabalho de \textcite{lewis2020retrieval}, para basear as respostas
estritamente nos documentos fornecidos e minimizar a ocorrência de
alucinações.

As principais funcionalidades incluem:

\begin{itemize}
\tightlist
\item
  \textbf{Resumos Multimodais (Áudio e Vídeo)}: Geração de conversas
  naturais entre especialistas no formato de \textbf{Resumo em Áudio}
  (estilo \emph{podcast}) e \textbf{Resumo em Vídeo}, discutindo os
  temas centrais do capítulo, como as diferenças entre processamento de
  imagens e visão computacional, ou a interpretação de transformações
  como a limiarização e o método de Otsu.
\item
  \textbf{Visualização de Estruturas (Mapa Mental e Infográfico)}:
  Criação automática de diagramas que conectam visualmente os conceitos,
  por exemplo, o fluxo de processamento desde a captura da imagem
  digital, passando pela conversão para tons de cinza, limiarização e
  segmentação binária.
\item
  \textbf{Ferramentas de Avaliação (Teste e Cartões Didáticos)}: Geração
  de \textbf{Testes} de múltipla escolha e \textbf{Cartões Didáticos}
  (\emph{flashcards}) para fixação de conhecimento, baseados no texto
  autoral (ex.: perguntas sobre a fórmula de conversão RGB→cinza ou
  sobre o funcionamento do limiar global e de Otsu).
\item
  \textbf{Apoio à Apresentação (\emph{Slides} e Relatórios)}: Auxílio na
  estruturação de \textbf{Apresentações de \emph{Slides}} e na redação
  de \textbf{Relatórios} técnicos e \textbf{\emph{Briefings}},
  facilitando a comunicação de resultados de experimentos com imagens.
\item
  \textbf{Análise de Dados (Tabela de Dados)}: Organização de dados
  extraídos do texto em tabelas estruturadas, auxiliando na compreensão
  de exemplos práticos, como a comparação entre diferentes valores de
  limiar.
\item
  \textbf{\emph{Chat} Contextualizado}: Permite o questionamento direto
  sobre o código e a teoria, como: \emph{``Como implementar a conversão
  de RGB para tons de cinza usando os pesos do padrão ITU‑R BT.601?''}
  ou \emph{``O que acontece com a imagem binária se eu escolher um
  limiar T=200 em vez de T=128?''}.
\end{itemize}

\begin{figure}

\centering{

\includegraphics[width=0.8\linewidth,height=\textheight,keepaspectratio]{cap01/imagens/fig-1-notebooklm.png}

}

\caption{\label{fig-notebooklm}Estúdio do NotebookLM.}

\end{figure}%

\section{Lista de Exercícios}\label{lista-de-exercuxedcios}

\begin{enumerate}
\def\labelenumi{\arabic{enumi}.}
\item
  (15\%) Com suas próprias palavras, defina \textbf{imagem digital} e
  \textbf{pixel}. Dê um exemplo concreto de como uma imagem colorida
  (RGB) é representada matricialmente no computador.
\item
  (15\%) Explique as diferenças entre \textbf{imagem binária},
  \textbf{tons de cinza (8 bits)} e \textbf{colorida RGB}, indicando a
  faixa de valores possíveis para cada pixel em cada tipo.
\item
  (20\%) Considerando a fórmula de conversão RGB → tons de cinza do
  padrão ITU‑R BT.601: \[g = 0.299\,R + 0.587\,G + 0.114\,B\] Calcule o
  valor do pixel em tons de cinza para \((R,G,B) = (80, 180, 30)\).
  Arredonde para o inteiro mais próximo.
\item
  (20\%) O que é \textbf{limiarização} (\emph{thresholding})? Explique a
  diferença entre escolher um limiar \(T\) fixo (ex.: \(T=128\)) e
  utilizar o \textbf{método de Otsu} para determinação automática do
  limiar. Em poucas palavras, como o método de Otsu escolhe o limiar?
\item
  (15\%) No contexto da biblioteca didática \texttt{mm} discutida no
  capítulo, responda:

  \begin{itemize}
  \tightlist
  \item
    \begin{enumerate}
    \def\labelenumii{\alph{enumii})}
    \tightlist
    \item
      (7,5\%) Como se acessa o valor do pixel na posição (linha=50,
      coluna=60) de uma imagem em tons de cinza \texttt{img\_gray}?
    \end{enumerate}
  \item
    \begin{enumerate}
    \def\labelenumii{\alph{enumii})}
    \setcounter{enumii}{1}
    \tightlist
    \item
      (7,5\%) Qual a vantagem de usar \texttt{mm.threshold(img\_gray)}
      sem passar o limiar? Compare com a chamada equivalente no OpenCV.
    \end{enumerate}
  \end{itemize}
\item
  (15\%) O que a propriedade \texttt{img.shape} retorna para uma imagem
  NumPy no formato RGB? Dê um exemplo concreto com uma imagem de 640×480
  pixels.
\end{enumerate}

\section{Referências do Capítulo}\label{referuxeancias-do-capuxedtulo}

Este capítulo baseou‑se principalmente nas seguintes obras de
processamento de imagens e visão computacional:

\begin{itemize}
\tightlist
\item
  \textcite{gonzalez2018digital} para fundamentos de
  \textbf{Processamento Digital de Imagens} (PDI).
\item
  \textcite{russ2023} para aspectos práticos de \textbf{processamento e
  análise de imagens} (The Image Processing Handbook, 7ª ed., 2016).
\item
  \textcite{szeliski2022} para \textbf{visão computacional} e algoritmos
  fundamentais (Computer Vision: Algorithms and Applications, 2ª ed.,
  2022).
\item
  \textcite{bradski2016} para a biblioteca \textbf{OpenCV} e exemplos
  práticos em Python (Learning OpenCV 4, 2ª ed., 2016).
\item
  \textcite{lewis2020retrieval} para o conceito de \textbf{geração
  aumentada por recuperação (RAG)}, aplicado ao NotebookLM.
\end{itemize}

\chapter{Da Captura ao Pixel -- Amostragem, Quantização e
Conectividade}\label{da-captura-ao-pixel-amostragem-quantizauxe7uxe3o-e-conectividade}

\href{https://colab.research.google.com/github/fzampirolli/pdi-vc/blob/master/notebooks_alunos/cap02/cap02_aluno.ipynb}{
\includegraphics[width=0.16\linewidth,height=\textheight,keepaspectratio]{cap02/imagens/colab-badge.png}}
\href{https://github.com/fzampirolli/pdi-vc}{
\includegraphics[width=0.2\linewidth,height=\textheight,keepaspectratio]{cap02/imagens/github-badge.png}}

\section{2.1. Objetivos}\label{objetivos-1}

Ao final deste capítulo, você será capaz de:

\begin{itemize}
\tightlist
\item
  Explicar o modelo simples de formação da imagem baseado em iluminação
  e refletância.
\item
  Diferenciar os principais elementos da percepção visual humana e dos
  sensores digitais.
\item
  Compreender os processos de \textbf{amostragem} (discretização do
  espaço) e \textbf{quantização} (discretização da intensidade).
\item
  Descrever as relações topológicas entre pixels: vizinhança,
  adjacência, conectividade e distâncias.
\item
  Realizar transformações geométricas básicas (translação, rotação,
  escala) preservando a qualidade da imagem.
\item
  Visualizar na prática os efeitos da variação da resolução espacial e
  da profundidade de bits.
\end{itemize}

\begin{center}\rule{0.5\linewidth}{0.5pt}\end{center}

\section{2.2. O Olho e a Câmera -- Elementos da Percepção
Visual}\label{o-olho-e-a-cuxe2mera-elementos-da-percepuxe7uxe3o-visual}

A formação de uma imagem digital começa com a captura da luz refletida
pelos objetos. Tanto o olho humano quanto as câmeras digitais seguem
princípios semelhantes, mas com diferenças importantes.

\subsection{2.2.1. Aspectos da visão
humana}\label{aspectos-da-visuxe3o-humana}

O olho é um sistema óptico complexo. A luz passa pela córnea, pelo humor
aquoso, pelo cristalino (que ajusta o foco) e pelo humor vítreo, até
atingir a retina. As células fotorreceptoras (cones e bastonetes)
convertem a luz em impulsos nervosos. Os cones são responsáveis pela
visão de cores (três tipos sensíveis ao vermelho, verde e azul) e os
bastonetes pela visão em baixa luminosidade (apenas intensidade). O
ponto cego é a região de saída do nervo óptico, sem fotorreceptores.

\subsection{2.2.2. Sensores digitais}\label{sensores-digitais}

As câmeras digitais utilizam sensores CCD ou CMOS que substituem a
retina. Cada elemento do sensor (fotossítio) acumula carga elétrica
proporcional à luz incidente. Um filtro de Bayer (matriz de filtros
coloridos) permite capturar as três componentes de cor. O processo de
leitura analógico-digital converte a carga em um valor numérico
(quantização). Diferentemente do olho, o sensor digital tem resolução
espacial fixa (número de pixels) e profundidade de bits definida (ex.: 8
bits por canal).

\begin{quote}
\textbf{Curiosidade:} O olho humano possui cerca de 120 milhões de
bastonetes e 6 milhões de cones, mas a resolução efetiva da visão
central (fóvea) é equivalente a aproximadamente 120\,×\,120 pixels de
alta definição. O processamento cerebral realiza enorme
pós-processamento.
\end{quote}

\begin{center}\rule{0.5\linewidth}{0.5pt}\end{center}

\section{2.3. O Modelo Matemático da Formação da
Imagem}\label{o-modelo-matemuxe1tico-da-formauxe7uxe3o-da-imagem}

Uma imagem pode ser modelada como o produto de duas funções:

{[} f(x,y) = i(x,y) \cdot r(x,y) {]}

onde:

\begin{itemize}
\tightlist
\item
  (i(x,y)) é a \textbf{iluminação} incidente sobre a cena (energia
  luminosa por unidade de área). Depende da fonte de luz.
\item
  (r(x,y)) é a \textbf{refletância} do objeto (fração da luz refletida
  em cada ponto). Depende das propriedades do material e da superfície.
\end{itemize}

Na prática, as duas componentes variam em faixas distintas: (i(x,y))
varia lentamente no espaço, enquanto (r(x,y)) pode variar rapidamente
(texturas, bordas). O processamento de imagens frequentemente tenta
separar ou compensar essas componentes (ex.: correção de iluminação não
uniforme).

\begin{center}\rule{0.5\linewidth}{0.5pt}\end{center}

\section{2.4. Digitalização: Amostragem e
Quantização}\label{digitalizauxe7uxe3o-amostragem-e-quantizauxe7uxe3o}

Para transformar uma cena contínua em uma imagem digital, dois processos
são necessários:

\subsection{2.4.1. Amostragem -- Discretização do
espaço}\label{amostragem-discretizauxe7uxe3o-do-espauxe7o}

A amostragem consiste em medir o valor da função (f(x,y)) em pontos
igualmente espaçados, formando uma matriz de (M) linhas e (N) colunas.
Cada elemento dessa matriz é um \textbf{pixel}. A \textbf{resolução
espacial} é dada por (M \times N). Quanto maior a resolução, mais
detalhes espaciais são preservados, mas maior também o custo
computacional e de armazenamento.

\textbf{Exemplo prático:} Uma imagem de 640\,×\,480 pixels tem resolução
espacial de 307\,200 pixels. Reduzir a resolução para 320\,×\,240
(subamostragem) reduz os detalhes e pode causar o efeito de
\emph{aliasing} (serrilhado). O aumento da resolução (sobre‐amostragem)
só é possível por interpolação, que não adiciona informação real.

\subsection{2.4.2. Quantização -- Discretização da
intensidade}\label{quantizauxe7uxe3o-discretizauxe7uxe3o-da-intensidade}

A quantização associa a cada pixel um valor numérico discreto,
geralmente representado por um inteiro de (b) bits. A
\textbf{profundidade de bits} ((b)) define o número de níveis de
intensidade: (2\^{}b) níveis. Imagens em tons de cinza costumam usar 8
bits (256 níveis). Imagens coloridas usam três canais de 8 bits (24 bits
no total).

\textbf{Ilustração:} Se usarmos apenas 1 bit por pixel (preto e branco),
perdemos todos os tons intermediários. Com 2 bits (4 níveis) já se
percebe degradês grosseiros. Com 8 bits, o olho humano dificilmente
percebe a discretização (visão contínua).

\begin{quote}
\textbf{Erro de quantização} é a diferença entre o valor analógico real
e o valor discreto atribuído. Ele se manifesta como ruído de
quantização, visível em regiões com gradiente suave quando se usa poucos
bits (efeito de ``posterização'').
\end{quote}

\begin{center}\rule{0.5\linewidth}{0.5pt}\end{center}

\section{2.5. Relações entre Pixels -- Topologia da
Imagem}\label{relauxe7uxf5es-entre-pixels-topologia-da-imagem}

Os pixels não são elementos isolados; suas posições relativas definem
importantes conceitos para processamento.

\subsection{2.5.1. Vizinhança}\label{vizinhanuxe7a}

Dado um pixel de coordenadas ((x,y)), definem‑se dois tipos principais
de vizinhança (para imagens em grid retangular):

\begin{itemize}
\tightlist
\item
  \textbf{Vizinhança‑4} (vizinhança de von Neumann): inclui os pixels
  nas posições ((x-1,y)), ((x+1,y)), ((x,y-1)), ((x,y+1)).
\item
  \textbf{Vizinhança‑8} (vizinhança de Moore): inclui todos os oito
  pixels adjacentes (acrescenta os quatro diagonais).
\end{itemize}

A escolha da vizinhança influencia operações como detecção de bordas,
cálculo de gradientes e conectividade.

\subsection{2.5.2. Adjacência, conectividade e
caminhos}\label{adjacuxeancia-conectividade-e-caminhos}

Dois pixels são \textbf{adjacentes} se estão em contato segundo uma
vizinhança definida e satisfazem um certo critério de valor (ex.: mesmo
nível de intensidade). Uma \textbf{conectividade} define uma relação de
equivalência entre pixels que formam uma região conexa. Um
\textbf{caminho} é uma sequência de pixels adjacentes. A conectividade‑4
e conectividade‑8 podem produzir resultados diferentes na segmentação e
no cálculo de componentes conexas (etiquetagem).

\subsection{2.5.3. Distâncias entre
pixels}\label{distuxe2ncias-entre-pixels}

Várias métricas podem ser usadas para medir distância:

\begin{itemize}
\tightlist
\item
  \textbf{Distância Euclidiana:} ( \sqrt{(x_1-x_2)^2 + (y_1-y_2)^2} ).
\item
  \textbf{Distância Manhattan (city block):} (
  \textbar x\_1-x\_2\textbar{} + \textbar y\_1-y\_2\textbar{} ).
\item
  \textbf{Distância Chebyshev (tabuleiro de xadrez):} (
  \max(\textbar x\_1-x\_2\textbar, \textbar y\_1-y\_2\textbar) ).
\end{itemize}

Essas distâncias são usadas em algoritmos de interpolação, transformada
de distância, crescimento de regiões, etc.

\begin{center}\rule{0.5\linewidth}{0.5pt}\end{center}

\section{2.6. Armazenamento de Imagens}\label{armazenamento-de-imagens}

Uma imagem digital pode ser armazenada em diversos formatos, cada um com
características próprias:

\begin{longtable}[]{@{}
  >{\raggedright\arraybackslash}p{(\linewidth - 4\tabcolsep) * \real{0.2368}}
  >{\raggedright\arraybackslash}p{(\linewidth - 4\tabcolsep) * \real{0.4211}}
  >{\raggedright\arraybackslash}p{(\linewidth - 4\tabcolsep) * \real{0.3421}}@{}}
\toprule\noalign{}
\begin{minipage}[b]{\linewidth}\raggedright
Formato
\end{minipage} & \begin{minipage}[b]{\linewidth}\raggedright
Características
\end{minipage} & \begin{minipage}[b]{\linewidth}\raggedright
Uso típico
\end{minipage} \\
\midrule\noalign{}
\endhead
\bottomrule\noalign{}
\endlastfoot
\textbf{BMP} & Não comprimido (ou compressão simples). & Windows,
aplicações legado. \\
\textbf{PNG} & Compressão sem perdas (lossless). & Web, imagens com
transparência. \\
\textbf{JPEG} & Compressão com perdas (lossy), boa para fotografias. &
Fotos, imagens médicas (uso moderado). \\
\textbf{TIFF} & Suporta múltiplas camadas e compressão variada. &
Editoração, arquivamento. \\
\textbf{RAW} & Dados brutos do sensor, sem processamento. & Fotografia
profissional. \\
\end{longtable}

Os metadados de uma imagem incluem: largura, altura, profundidade de
bits, tipo de codificação de cor (grayscale, RGB, CMYK), informações de
calibração (ex.: resolução DPI), e metadados EXIF (data, modelo da
câmera, configurações). Ao ler uma imagem com a função
\texttt{mm.read()}, a biblioteca \texttt{morph} preserva automaticamente
os metadados relevantes.

\begin{center}\rule{0.5\linewidth}{0.5pt}\end{center}

\section{2.7. Transformações Geométricas
Básicas}\label{transformauxe7uxf5es-geomuxe9tricas-buxe1sicas}

As transformações geométricas alteram a posição dos pixels, mantendo os
valores de intensidade. São fundamentais para alinhamento, correção de
distorções e aumento de dados.

\subsection{2.7.1. Translação}\label{translauxe7uxe3o}

Desloca a imagem em (t\_x) e (t\_y) pixels: {[}

\begin{bmatrix} x' \\ y' \end{bmatrix}

=

\begin{bmatrix} 1 & 0 & t_x \\ 0 & 1 & t_y \end{bmatrix} \begin{bmatrix} x \\ y \\ 1 \end{bmatrix}

. {]} Na prática, cria‑se uma nova imagem maior e copia‑se os pixels
deslocados. Pixels que saem da área original são descartados e áreas
vazias são preenchidas (geralmente com 0 -- preto).

\subsection{2.7.2. Rotação}\label{rotauxe7uxe3o}

Gira a imagem em torno da origem (ou de seu centro) por um ângulo
(\theta). A transformação direta pode produzir buracos; na prática,
usa‑se a transformação inversa (mapeamento de cada pixel da imagem
destino para a origem) e interpolação (vizinho mais próximo, bilinear,
bicúbica). A rotação geralmente requer redimensionamento da imagem para
que nenhuma área seja cortada.

\subsection{2.7.3. Escala
(redimensionamento)}\label{escala-redimensionamento}

Altera o tamanho da imagem pelo fator (s\_x) (horizontal) e (s\_y)
(vertical). O redimensionamento exige interpolação:

\begin{itemize}
\tightlist
\item
  \textbf{Vizinho mais próximo:} rápido, mas produz efeito de blocos.
\item
  \textbf{Bilinear:} suaviza melhor, é um bom compromisso.
\item
  \textbf{Bicúbica:} mais suave, usada em softwares profissionais.
\end{itemize}

A biblioteca \texttt{morph} oferece funções como
\texttt{mm.resize(img,\ scale,\ interpolation)} e
\texttt{mm.rotate(img,\ angle,\ center,\ interp)}.

\begin{center}\rule{0.5\linewidth}{0.5pt}\end{center}

\section{2.8. Laboratório Prático: Visualizando Efeitos da Amostragem e
Quantização}\label{laboratuxf3rio-pruxe1tico-visualizando-efeitos-da-amostragem-e-quantizauxe7uxe3o}

Neste experimento, vamos explorar como a redução da resolução espacial e
da profundidade de bits afetam a qualidade visual.

\subsection{2.8.1. Código para subamostragem (redução de
resolução)}\label{cuxf3digo-para-subamostragem-reduuxe7uxe3o-de-resoluuxe7uxe3o}

\begin{Shaded}
\begin{Highlighting}[]
\ImportTok{import}\NormalTok{ numpy }\ImportTok{as}\NormalTok{ np}
\ImportTok{import}\NormalTok{ matplotlib.pyplot }\ImportTok{as}\NormalTok{ plt}
\ImportTok{import}\NormalTok{ morph}

\CommentTok{\# Carregar uma imagem (ex.: Lenna)}
\NormalTok{img }\OperatorTok{=}\NormalTok{ mm.read(}\StringTok{\textquotesingle{}https://upload.wikimedia.org/wikipedia/en/7/7d/Lenna\_\%28test\_image\%29.png\textquotesingle{}}\NormalTok{)}
\NormalTok{img\_gray }\OperatorTok{=}\NormalTok{ mm.rgb2gray(img)}

\CommentTok{\# Função para subamostrar por fator (ex.: 2, 4, 8)}
\KeywordTok{def}\NormalTok{ subsample(image, factor):}
\NormalTok{    h, w }\OperatorTok{=}\NormalTok{ image.shape}
\NormalTok{    new\_h, new\_w }\OperatorTok{=}\NormalTok{ h }\OperatorTok{//}\NormalTok{ factor, w }\OperatorTok{//}\NormalTok{ factor}
    \CommentTok{\# Criar imagem reduzida: pegar um pixel a cada \textquotesingle{}factor\textquotesingle{}}
\NormalTok{    reduced }\OperatorTok{=}\NormalTok{ image[::factor, ::factor]}
    \CommentTok{\# Redimensionar de volta ao tamanho original para comparação (interpolação vizinho mais próximo)}
    \ControlFlowTok{return}\NormalTok{ mm.resize(reduced, (h, w), method}\OperatorTok{=}\StringTok{\textquotesingle{}nearest\textquotesingle{}}\NormalTok{)}

\NormalTok{plt.figure(figsize}\OperatorTok{=}\NormalTok{(}\DecValTok{15}\NormalTok{, }\DecValTok{5}\NormalTok{))}
\NormalTok{factors }\OperatorTok{=}\NormalTok{ [}\DecValTok{1}\NormalTok{, }\DecValTok{2}\NormalTok{, }\DecValTok{4}\NormalTok{, }\DecValTok{8}\NormalTok{]}
\ControlFlowTok{for}\NormalTok{ i, f }\KeywordTok{in} \BuiltInTok{enumerate}\NormalTok{(factors):}
\NormalTok{    sub }\OperatorTok{=}\NormalTok{ subsample(img\_gray, f)}
\NormalTok{    plt.subplot(}\DecValTok{1}\NormalTok{, }\DecValTok{4}\NormalTok{, i}\OperatorTok{+}\DecValTok{1}\NormalTok{)}
\NormalTok{    mm.imshow(sub)}
\NormalTok{    plt.title(}\SpecialStringTok{f\textquotesingle{}Fator = }\SpecialCharTok{\{}\NormalTok{f}\SpecialCharTok{\}}\SpecialStringTok{\textquotesingle{}}\NormalTok{)}
\NormalTok{plt.show()}
\end{Highlighting}
\end{Shaded}

\subsection{2.8.2. Código para redução da profundidade de bits
(quantização)}\label{cuxf3digo-para-reduuxe7uxe3o-da-profundidade-de-bits-quantizauxe7uxe3o}

\begin{Shaded}
\begin{Highlighting}[]
\KeywordTok{def}\NormalTok{ quantize(image, bits):}
    \CommentTok{"""Reduz a imagem para 2\^{}bits níveis de cinza."""}
\NormalTok{    levels }\OperatorTok{=} \DecValTok{2} \OperatorTok{**}\NormalTok{ bits}
    \CommentTok{\# Normalizar, quantizar e denormalizar}
    \ControlFlowTok{return}\NormalTok{ np.floor(image }\OperatorTok{/} \DecValTok{256} \OperatorTok{*}\NormalTok{ levels) }\OperatorTok{/}\NormalTok{ levels }\OperatorTok{*} \DecValTok{255}

\NormalTok{plt.figure(figsize}\OperatorTok{=}\NormalTok{(}\DecValTok{15}\NormalTok{, }\DecValTok{5}\NormalTok{))}
\NormalTok{bits }\OperatorTok{=}\NormalTok{ [}\DecValTok{8}\NormalTok{, }\DecValTok{4}\NormalTok{, }\DecValTok{2}\NormalTok{, }\DecValTok{1}\NormalTok{]}
\ControlFlowTok{for}\NormalTok{ i, b }\KeywordTok{in} \BuiltInTok{enumerate}\NormalTok{(bits):}
\NormalTok{    quant }\OperatorTok{=}\NormalTok{ quantize(img\_gray, b)}
\NormalTok{    plt.subplot(}\DecValTok{1}\NormalTok{, }\DecValTok{4}\NormalTok{, i}\OperatorTok{+}\DecValTok{1}\NormalTok{)}
\NormalTok{    mm.imshow(quant)}
\NormalTok{    plt.title(}\SpecialStringTok{f\textquotesingle{}}\SpecialCharTok{\{}\NormalTok{b}\SpecialCharTok{\}}\SpecialStringTok{ bits (}\SpecialCharTok{\{}\DecValTok{2}\OperatorTok{**}\NormalTok{b}\SpecialCharTok{\}}\SpecialStringTok{ níveis)\textquotesingle{}}\NormalTok{)}
\NormalTok{plt.show()}
\end{Highlighting}
\end{Shaded}

\subsection{2.8.3. Discussão dos
resultados}\label{discussuxe3o-dos-resultados}

\begin{itemize}
\tightlist
\item
  \textbf{Subamostragem (fator 8):} observa‑se o efeito de
  \emph{aliasing} -- bordas serrilhadas e perda de detalhes finos. Em
  fator 4 já se nota blocos, mas a imagem ainda é reconhecível.
\item
  \textbf{Quantização com 1 bit:} a imagem vira preto e branco, perde‑se
  totalmente a graduação de tons. Com 2 bits (4 níveis) ainda se percebe
  um efeito de ``posters'', mas a silhueta geral se mantém. Com 4 bits
  (16 níveis) a qualidade é razoável; o olho humano dificilmente
  distingue 8 bits de 16 bits em imagens comuns.
\end{itemize}

Esses experimentos mostram o compromisso entre qualidade e custo de
armazenamento, essencial para sistemas de processamento e compressão.

\begin{center}\rule{0.5\linewidth}{0.5pt}\end{center}

\section{2.9. Exercícios Propostos}\label{exercuxedcios-propostos}

\begin{enumerate}
\def\labelenumi{\arabic{enumi}.}
\tightlist
\item
  \textbf{Amostragem} -- Carregue uma imagem, aplique subamostragem com
  fatores 2, 3 e 5. Compare visualmente e calcule o fator de redução no
  número de pixels.
\item
  \textbf{Quantização} -- Implemente uma função que reduza a imagem a 3
  bits (8 níveis). Visualize e descreva o ruído introduzido.
\item
  \textbf{Conectividade} -- Escreva um código que identifique
  componentes conexas em uma imagem binária usando vizinhança‑4 e
  vizinhança‑8. Explique as diferenças.
\item
  \textbf{Transformação geométrica} -- Aplique rotação de 30° e 45° em
  uma imagem, utilizando interpolação bilinear. Compare com a rotação
  sem interpolação (vizinho mais próximo) e comente sobre a suavidade.
\end{enumerate}

\begin{center}\rule{0.5\linewidth}{0.5pt}\end{center}

\section{2.10. Referências}\label{referuxeancias}

\begin{itemize}
\tightlist
\item
  GONZALEZ, R. C.; WOODS, R. E. \emph{Processamento Digital de Imagens}.
  3. ed.~Pearson, 2010.
\item
  SZELISKI, R. \emph{Computer Vision: Algorithms and Applications}.
  Springer, 2010.
\item
  Otsu, N. (1979). \emph{A threshold selection method from gray‑level
  histograms}. IEEE Trans. Syst. Man Cybern., 9(1), 62--66. (citado no
  Capítulo 1).
\end{itemize}

\begin{center}\rule{0.5\linewidth}{0.5pt}\end{center}

\textbf{Nota:} As operações de convolução, filtragem e morfologia
matemática (erosão, dilatação) serão tratadas no \textbf{Capítulo 3 --
Operações Espaciais em Imagens}. O presente capítulo concentrou‑se na
formação, digitalização e topologia básica da imagem.

\section{Objetivos}\label{objetivos-2}

Ao final deste capítulo, o aluno será capaz de:

\begin{itemize}
\tightlist
\item
  Compreender a representação matemática de imagens \(f(x,y)\)
\item
  Ler, exibir e salvar imagens
\item
  Trabalhar com imagens binárias, tons de cinza e coloridas (RGB)
\item
  Acessar e modificar pixels individualmente
\end{itemize}

\section{O que é uma Imagem
Digital?}\label{o-que-uxe9-uma-imagem-digital-1}

Uma imagem digital é uma função discreta:

\begin{equation}\phantomsection\label{eq-1-imagem}{f: \mathbb{Z}^2 \rightarrow \mathbb{Z}^k
}\end{equation}

onde \((x,y)\) são coordenadas espaciais e \(f(x,y)\) é a intensidade do
pixel. Para imagens em \textbf{tons de cinza} \(k=1\) com valores em
\([0, 255]\); para imagens \textbf{coloridas} (RGB) \(k=3\), ver
Figura~\ref{fig-1-representacao}.

\begin{figure}

\centering{

\includegraphics[width=0.7\linewidth,height=\textheight,keepaspectratio]{cap02/imagens/lena.jpg}

}

\caption{\label{fig-1-representacao}Exemplo de Transformação Não Linear
do Espaço de Entrada para o Espaço de Características.}

\end{figure}%

\section{Configuração do
Ambiente}\label{configurauxe7uxe3o-do-ambiente-1}

\begin{Shaded}
\begin{Highlighting}[]
\CommentTok{\# Instalação (Google Colab — executar uma vez)}
\CommentTok{\# !pip install opencv{-}python{-}headless scikit{-}image matplotlib numpy}

\ImportTok{import}\NormalTok{ sys}
\ImportTok{import}\NormalTok{ importlib}
\ImportTok{import}\NormalTok{ numpy }\ImportTok{as}\NormalTok{ np}
\ImportTok{import}\NormalTok{ cv2}
\ImportTok{import}\NormalTok{ matplotlib.pyplot }\ImportTok{as}\NormalTok{ plt}

\NormalTok{sys.path.insert(}\DecValTok{0}\NormalTok{, }\StringTok{\textquotesingle{}../../morph\textquotesingle{}}\NormalTok{)}
\ImportTok{import}\NormalTok{ morph}
\NormalTok{importlib.}\BuiltInTok{reload}\NormalTok{(morph)}
\ImportTok{from}\NormalTok{ morph }\ImportTok{import}\NormalTok{ mm}

\BuiltInTok{print}\NormalTok{(}\SpecialStringTok{f\textquotesingle{}OpenCV }\SpecialCharTok{\{}\NormalTok{cv2}\SpecialCharTok{.}\NormalTok{\_\_version\_\_}\SpecialCharTok{\}}\SpecialStringTok{ | NumPy }\SpecialCharTok{\{}\NormalTok{np}\SpecialCharTok{.}\NormalTok{\_\_version\_\_}\SpecialCharTok{\}}\SpecialStringTok{\textquotesingle{}}\NormalTok{)}
\BuiltInTok{print}\NormalTok{(}\StringTok{"Carregando imagem..."}\NormalTok{)}
\BuiltInTok{print}\NormalTok{(}\BuiltInTok{help}\NormalTok{(mm.show))}
\end{Highlighting}
\end{Shaded}

\begin{verbatim}
OpenCV 4.12.0 | NumPy 2.2.6
Carregando imagem...
Help on function show in module morph:

show(*args, title=None)
    This function will draw images f
    input: <*args> set of images f_i, where i>0 is binary image
          <title> optional string title for the plot
    output: image drawing
    Example:
    f1, f2 = np.zeros((100, 100,3)),  np.zeros((100, 100))
    f2[50:60, 50:60] = 1
    mm.show(f1, f2, title='Exemplo')

None
\end{verbatim}

\section{1.3 Lendo e Exibindo Imagens}\label{lendo-e-exibindo-imagens-1}

A operação mais básica em PDI é a leitura de uma imagem. A função
\texttt{mm.read()} da \texttt{morph.py} aceita caminhos locais e URLs.

\begin{Shaded}
\begin{Highlighting}[]
\ImportTok{import}\NormalTok{ urllib.request, os}

\CommentTok{\# Baixar imagem de exemplo}
\NormalTok{url }\OperatorTok{=} \StringTok{\textquotesingle{}https://raw.githubusercontent.com/opencv/opencv/master/samples/data/lena.jpg\textquotesingle{}}
\NormalTok{os.makedirs(}\StringTok{\textquotesingle{}imagens\textquotesingle{}}\NormalTok{, exist\_ok}\OperatorTok{=}\VariableTok{True}\NormalTok{)}
\ControlFlowTok{if} \KeywordTok{not}\NormalTok{ os.path.exists(}\StringTok{\textquotesingle{}imagens/lena.jpg\textquotesingle{}}\NormalTok{):}
\NormalTok{    urllib.request.urlretrieve(url, }\StringTok{\textquotesingle{}imagens/lena.jpg\textquotesingle{}}\NormalTok{)}

\CommentTok{\# Ler imagem com morph.py}
\NormalTok{img }\OperatorTok{=}\NormalTok{ mm.read(}\StringTok{\textquotesingle{}imagens/lena.jpg\textquotesingle{}}\NormalTok{)}

\BuiltInTok{print}\NormalTok{(}\SpecialStringTok{f\textquotesingle{}Shape  : }\SpecialCharTok{\{}\NormalTok{img}\SpecialCharTok{.}\NormalTok{shape}\SpecialCharTok{\}}\SpecialStringTok{\textquotesingle{}}\NormalTok{)    }\CommentTok{\# (altura, largura, canais)}
\BuiltInTok{print}\NormalTok{(}\SpecialStringTok{f\textquotesingle{}Dtype  : }\SpecialCharTok{\{}\NormalTok{img}\SpecialCharTok{.}\NormalTok{dtype}\SpecialCharTok{\}}\SpecialStringTok{\textquotesingle{}}\NormalTok{)    }\CommentTok{\# uint8}
\BuiltInTok{print}\NormalTok{(}\SpecialStringTok{f\textquotesingle{}Min/Max: }\SpecialCharTok{\{}\NormalTok{img}\SpecialCharTok{.}\BuiltInTok{min}\NormalTok{()}\SpecialCharTok{\}}\SpecialStringTok{ / }\SpecialCharTok{\{}\NormalTok{img}\SpecialCharTok{.}\BuiltInTok{max}\NormalTok{()}\SpecialCharTok{\}}\SpecialStringTok{\textquotesingle{}}\NormalTok{)}

\NormalTok{img\_rgb }\OperatorTok{=}\NormalTok{ cv2.cvtColor(img, cv2.COLOR\_BGR2RGB)}
\NormalTok{mm.show(img\_rgb, title}\OperatorTok{=}\StringTok{\textquotesingle{}Imagem original\textquotesingle{}}\NormalTok{)}
\end{Highlighting}
\end{Shaded}

\begin{verbatim}
Shape  : (512, 512, 3)
Dtype  : uint8
Min/Max: 0 / 255
\end{verbatim}

\begin{figure}[H]

\centering{

\pandocbounded{\includegraphics[keepaspectratio]{cap02/cap02.py.pt_files/figure-pdf/fig-1-lena-output-2.png}}

}

\caption{\label{fig-1-lena}MVS Linear --- Margem Máxima entre Íris
Setosa e Versicolor (atributos de pétala).}

\end{figure}%

\section{1.4 Tipos de Imagem}\label{tipos-de-imagem}

\subsection{1.4.1 Conversão para Tons de
Cinza}\label{conversuxe3o-para-tons-de-cinza-1}

A conversão de RGB para tons de cinza aplica pesos perceptuais definidos
pela ITU-R BT.601:

\begin{equation}\phantomsection\label{eq-1-cinza}{g = 0.299\,R + 0.587\,G + 0.114\,B}\end{equation}

\begin{Shaded}
\begin{Highlighting}[]
\CommentTok{\# Converter para tons de cinza}
\NormalTok{img\_gray }\OperatorTok{=}\NormalTok{ mm.gray(img)}
\BuiltInTok{print}\NormalTok{(}\SpecialStringTok{f\textquotesingle{}Shape cinza: }\SpecialCharTok{\{}\NormalTok{img\_gray}\SpecialCharTok{.}\NormalTok{shape}\SpecialCharTok{\}}\SpecialStringTok{\textquotesingle{}}\NormalTok{)}

\NormalTok{mm.show([img, img\_gray], title}\OperatorTok{=}\NormalTok{[}\StringTok{\textquotesingle{}Original (RGB)\textquotesingle{}}\NormalTok{, }\StringTok{\textquotesingle{}Tons de Cinza\textquotesingle{}}\NormalTok{])}
\end{Highlighting}
\end{Shaded}

\subsection{1.4.2 Limiarização
(Binarização)}\label{limiarizauxe7uxe3o-binarizauxe7uxe3o}

A limiarização converte uma imagem em tons de cinza para binária:

\begin{equation}\phantomsection\label{eq-1-limiar}{g(x,y) = \begin{cases} 255 & \text{se } f(x,y) \geq T \\ 0 & \text{caso contrário} \end{cases}}\end{equation}

O método de \textbf{Otsu} \autocite{otsu1979} determina \(T\)
automaticamente maximizando a variância inter-classes do histograma.

\begin{Shaded}
\begin{Highlighting}[]
\CommentTok{\# Limiarização com T fixo}
\NormalTok{T }\OperatorTok{=} \DecValTok{128}
\NormalTok{img\_bin }\OperatorTok{=}\NormalTok{ mm.threshold(img\_gray, T)}

\CommentTok{\# Limiarização Otsu (T automático)}
\NormalTok{\_, img\_otsu }\OperatorTok{=}\NormalTok{ cv2.threshold(img\_gray, }\DecValTok{0}\NormalTok{, }\DecValTok{255}\NormalTok{,}
\NormalTok{                             cv2.THRESH\_BINARY }\OperatorTok{+}\NormalTok{ cv2.THRESH\_OTSU)}

\BuiltInTok{print}\NormalTok{(}\SpecialStringTok{f\textquotesingle{}Valores únicos (T=}\SpecialCharTok{\{}\NormalTok{T}\SpecialCharTok{\}}\SpecialStringTok{): }\SpecialCharTok{\{}\BuiltInTok{sorted}\NormalTok{(}\BuiltInTok{set}\NormalTok{(img\_bin.flatten()))}\SpecialCharTok{\}}\SpecialStringTok{\textquotesingle{}}\NormalTok{)}

\NormalTok{mm.show([img\_gray, img\_bin, img\_otsu],}
\NormalTok{        title}\OperatorTok{=}\NormalTok{[}\SpecialStringTok{f\textquotesingle{}Cinza\textquotesingle{}}\NormalTok{, }\SpecialStringTok{f\textquotesingle{}Binária (T=}\SpecialCharTok{\{}\NormalTok{T}\SpecialCharTok{\}}\SpecialStringTok{)\textquotesingle{}}\NormalTok{, }\StringTok{\textquotesingle{}Otsu\textquotesingle{}}\NormalTok{])}
\end{Highlighting}
\end{Shaded}

\section{1.5 Acesso a Pixels}\label{acesso-a-pixels-1}

Em Python/NumPy, o acesso a pixels usa indexação de arrays
\texttt{img{[}linha,\ coluna{]}}.

\begin{Shaded}
\begin{Highlighting}[]
\NormalTok{r, c }\OperatorTok{=} \DecValTok{100}\NormalTok{, }\DecValTok{100}

\CommentTok{\# Acessar pixel}
\NormalTok{pixel\_cinza }\OperatorTok{=}\NormalTok{ img\_gray[r, c]}
\NormalTok{pixel\_rgb   }\OperatorTok{=}\NormalTok{ img[r, c]       }\CommentTok{\# array [R, G, B]}

\BuiltInTok{print}\NormalTok{(}\SpecialStringTok{f\textquotesingle{}Pixel cinza (}\SpecialCharTok{\{}\NormalTok{r}\SpecialCharTok{\}}\SpecialStringTok{,}\SpecialCharTok{\{}\NormalTok{c}\SpecialCharTok{\}}\SpecialStringTok{): }\SpecialCharTok{\{}\NormalTok{pixel\_cinza}\SpecialCharTok{\}}\SpecialStringTok{\textquotesingle{}}\NormalTok{)}
\BuiltInTok{print}\NormalTok{(}\SpecialStringTok{f\textquotesingle{}Pixel RGB   (}\SpecialCharTok{\{}\NormalTok{r}\SpecialCharTok{\}}\SpecialStringTok{,}\SpecialCharTok{\{}\NormalTok{c}\SpecialCharTok{\}}\SpecialStringTok{): R=}\SpecialCharTok{\{}\NormalTok{pixel\_rgb[}\DecValTok{0}\NormalTok{]}\SpecialCharTok{\}}\SpecialStringTok{  G=}\SpecialCharTok{\{}\NormalTok{pixel\_rgb[}\DecValTok{1}\NormalTok{]}\SpecialCharTok{\}}\SpecialStringTok{  B=}\SpecialCharTok{\{}\NormalTok{pixel\_rgb[}\DecValTok{2}\NormalTok{]}\SpecialCharTok{\}}\SpecialStringTok{\textquotesingle{}}\NormalTok{)}

\CommentTok{\# Imagem sintética 5×5 com pixel central branco}
\NormalTok{syn }\OperatorTok{=}\NormalTok{ np.zeros((}\DecValTok{5}\NormalTok{, }\DecValTok{5}\NormalTok{), dtype}\OperatorTok{=}\NormalTok{np.uint8)}
\NormalTok{syn[}\DecValTok{2}\NormalTok{, }\DecValTok{2}\NormalTok{] }\OperatorTok{=} \DecValTok{255}
\BuiltInTok{print}\NormalTok{(}\StringTok{\textquotesingle{}}\CharTok{\textbackslash{}n}\StringTok{Matriz 5×5:\textquotesingle{}}\NormalTok{)}
\BuiltInTok{print}\NormalTok{(syn)}
\end{Highlighting}
\end{Shaded}

\section{1.6 Resumo}\label{resumo-1}

Neste capítulo foram apresentados os fundamentos de representação de
imagens digitais:

\begin{itemize}
\tightlist
\item
  Imagem digital = função \(f(x,y)\) mapeando coordenadas para
  intensidades
\item
  Tipos: binária (\(k=1\), \(\{0,255\}\)), tons de cinza (\(k=1\),
  \([0,255]\)), RGB (\(k=3\))
\item
  \texttt{morph.py}: biblioteca didática que simplifica operações de PDI
  em Python
\item
  Limiarização: operação fundamental de binarização
\end{itemize}

O Capítulo 2 estuda \textbf{histogramas e equalização}.

\subsection{Questões de Revisão}\label{questuxf5es-de-revisuxe3o}

\begin{enumerate}
\def\labelenumi{\arabic{enumi}.}
\tightlist
\item
  Qual a diferença entre imagem em tons de cinza e imagem binária?
\item
  Como o método de Otsu determina o limiar \(T\) automaticamente?
\item
  O que representa \texttt{img.shape} para uma imagem NumPy RGB?
\end{enumerate}

\printbibliography

\chapter{Operações Espaciais: Transformações Geométricas, Convolução e
Morfologia}\label{operauxe7uxf5es-espaciais-transformauxe7uxf5es-geomuxe9tricas-convoluuxe7uxe3o-e-morfologia}

\href{https://colab.research.google.com/github/fzampirolli/pdi-vc/blob/master/notebooks_alunos/cap03/cap03_aluno.ipynb}{
\includegraphics[width=0.16\linewidth,height=\textheight,keepaspectratio]{cap03/imagens/colab-badge.png}}
\href{https://github.com/fzampirolli/pdi-vc}{
\includegraphics[width=0.2\linewidth,height=\textheight,keepaspectratio]{cap03/imagens/github-badge.png}}

Objetivos: Aplicar transformações geométricas (rotação, escala,
translação); realizar operações de convolução e filtragem; introduzir
operações morfológicas (erosão, dilatação, abertura, fechamento).

Conteúdo:

Transformações geométricas (interpolação)

Convolução e filtros (média, gaussiano, bordas)

Morfologia matemática (erosão, dilatação, operações derivadas)

\section{Objetivos}\label{objetivos-3}

Ao final deste capítulo, o aluno será capaz de:

\begin{itemize}
\tightlist
\item
  Compreender a representação matemática de imagens \(f(x,y)\)
\item
  Ler, exibir e salvar imagens
\item
  Trabalhar com imagens binárias, tons de cinza e coloridas (RGB)
\item
  Acessar e modificar pixels individualmente
\end{itemize}

\section{O que é uma Imagem
Digital?}\label{o-que-uxe9-uma-imagem-digital-2}

Uma imagem digital é uma função discreta:

\begin{equation}\phantomsection\label{eq-1-imagem}{f: \mathbb{Z}^2 \rightarrow \mathbb{Z}^k
}\end{equation}

onde \((x,y)\) são coordenadas espaciais e \(f(x,y)\) é a intensidade do
pixel. Para imagens em \textbf{tons de cinza} \(k=1\) com valores em
\([0, 255]\); para imagens \textbf{coloridas} (RGB) \(k=3\), ver
Figura~\ref{fig-1-representacao}.

\begin{figure}

\centering{

\includegraphics[width=0.7\linewidth,height=\textheight,keepaspectratio]{cap03/imagens/lena.jpg}

}

\caption{\label{fig-1-representacao}Exemplo de Transformação Não Linear
do Espaço de Entrada para o Espaço de Características.}

\end{figure}%

\section{Configuração do
Ambiente}\label{configurauxe7uxe3o-do-ambiente-2}

\begin{Shaded}
\begin{Highlighting}[]
\CommentTok{\# Instalação (Google Colab — executar uma vez)}
\CommentTok{\# !pip install opencv{-}python{-}headless scikit{-}image matplotlib numpy}

\ImportTok{import}\NormalTok{ sys}
\ImportTok{import}\NormalTok{ importlib}
\ImportTok{import}\NormalTok{ numpy }\ImportTok{as}\NormalTok{ np}
\ImportTok{import}\NormalTok{ cv2}
\ImportTok{import}\NormalTok{ matplotlib.pyplot }\ImportTok{as}\NormalTok{ plt}

\NormalTok{sys.path.insert(}\DecValTok{0}\NormalTok{, }\StringTok{\textquotesingle{}../../morph\textquotesingle{}}\NormalTok{)}
\ImportTok{import}\NormalTok{ morph}
\NormalTok{importlib.}\BuiltInTok{reload}\NormalTok{(morph)}
\ImportTok{from}\NormalTok{ morph }\ImportTok{import}\NormalTok{ mm}

\BuiltInTok{print}\NormalTok{(}\SpecialStringTok{f\textquotesingle{}OpenCV }\SpecialCharTok{\{}\NormalTok{cv2}\SpecialCharTok{.}\NormalTok{\_\_version\_\_}\SpecialCharTok{\}}\SpecialStringTok{ | NumPy }\SpecialCharTok{\{}\NormalTok{np}\SpecialCharTok{.}\NormalTok{\_\_version\_\_}\SpecialCharTok{\}}\SpecialStringTok{\textquotesingle{}}\NormalTok{)}
\BuiltInTok{print}\NormalTok{(}\StringTok{"Carregando imagem..."}\NormalTok{)}
\BuiltInTok{print}\NormalTok{(}\BuiltInTok{help}\NormalTok{(mm.show))}
\end{Highlighting}
\end{Shaded}

\begin{verbatim}
OpenCV 4.12.0 | NumPy 2.2.6
Carregando imagem...
Help on function show in module morph:

show(*args, title=None)
    This function will draw images f
    input: <*args> set of images f_i, where i>0 is binary image
          <title> optional string title for the plot
    output: image drawing
    Example:
    f1, f2 = np.zeros((100, 100,3)),  np.zeros((100, 100))
    f2[50:60, 50:60] = 1
    mm.show(f1, f2, title='Exemplo')

None
\end{verbatim}

\section{1.3 Lendo e Exibindo Imagens}\label{lendo-e-exibindo-imagens-2}

A operação mais básica em PDI é a leitura de uma imagem. A função
\texttt{mm.read()} da \texttt{morph.py} aceita caminhos locais e URLs.

\begin{Shaded}
\begin{Highlighting}[]
\ImportTok{import}\NormalTok{ urllib.request, os}

\CommentTok{\# Baixar imagem de exemplo}
\NormalTok{url }\OperatorTok{=} \StringTok{\textquotesingle{}https://raw.githubusercontent.com/opencv/opencv/master/samples/data/lena.jpg\textquotesingle{}}
\NormalTok{os.makedirs(}\StringTok{\textquotesingle{}imagens\textquotesingle{}}\NormalTok{, exist\_ok}\OperatorTok{=}\VariableTok{True}\NormalTok{)}
\ControlFlowTok{if} \KeywordTok{not}\NormalTok{ os.path.exists(}\StringTok{\textquotesingle{}imagens/lena.jpg\textquotesingle{}}\NormalTok{):}
\NormalTok{    urllib.request.urlretrieve(url, }\StringTok{\textquotesingle{}imagens/lena.jpg\textquotesingle{}}\NormalTok{)}

\CommentTok{\# Ler imagem com morph.py}
\NormalTok{img }\OperatorTok{=}\NormalTok{ mm.read(}\StringTok{\textquotesingle{}imagens/lena.jpg\textquotesingle{}}\NormalTok{)}

\BuiltInTok{print}\NormalTok{(}\SpecialStringTok{f\textquotesingle{}Shape  : }\SpecialCharTok{\{}\NormalTok{img}\SpecialCharTok{.}\NormalTok{shape}\SpecialCharTok{\}}\SpecialStringTok{\textquotesingle{}}\NormalTok{)    }\CommentTok{\# (altura, largura, canais)}
\BuiltInTok{print}\NormalTok{(}\SpecialStringTok{f\textquotesingle{}Dtype  : }\SpecialCharTok{\{}\NormalTok{img}\SpecialCharTok{.}\NormalTok{dtype}\SpecialCharTok{\}}\SpecialStringTok{\textquotesingle{}}\NormalTok{)    }\CommentTok{\# uint8}
\BuiltInTok{print}\NormalTok{(}\SpecialStringTok{f\textquotesingle{}Min/Max: }\SpecialCharTok{\{}\NormalTok{img}\SpecialCharTok{.}\BuiltInTok{min}\NormalTok{()}\SpecialCharTok{\}}\SpecialStringTok{ / }\SpecialCharTok{\{}\NormalTok{img}\SpecialCharTok{.}\BuiltInTok{max}\NormalTok{()}\SpecialCharTok{\}}\SpecialStringTok{\textquotesingle{}}\NormalTok{)}

\NormalTok{img\_rgb }\OperatorTok{=}\NormalTok{ cv2.cvtColor(img, cv2.COLOR\_BGR2RGB)}
\NormalTok{mm.show(img\_rgb, title}\OperatorTok{=}\StringTok{\textquotesingle{}Imagem original\textquotesingle{}}\NormalTok{)}
\end{Highlighting}
\end{Shaded}

\begin{verbatim}
Shape  : (512, 512, 3)
Dtype  : uint8
Min/Max: 0 / 255
\end{verbatim}

\begin{figure}[H]

\centering{

\pandocbounded{\includegraphics[keepaspectratio]{cap03/cap03.py.pt_files/figure-pdf/fig-1-lena-output-2.png}}

}

\caption{\label{fig-1-lena}MVS Linear --- Margem Máxima entre Íris
Setosa e Versicolor (atributos de pétala).}

\end{figure}%

\section{1.4 Tipos de Imagem}\label{tipos-de-imagem-1}

\subsection{1.4.1 Conversão para Tons de
Cinza}\label{conversuxe3o-para-tons-de-cinza-2}

A conversão de RGB para tons de cinza aplica pesos perceptuais definidos
pela ITU-R BT.601:

\begin{equation}\phantomsection\label{eq-1-cinza}{g = 0.299\,R + 0.587\,G + 0.114\,B}\end{equation}

\begin{Shaded}
\begin{Highlighting}[]
\CommentTok{\# Converter para tons de cinza}
\NormalTok{img\_gray }\OperatorTok{=}\NormalTok{ mm.gray(img)}
\BuiltInTok{print}\NormalTok{(}\SpecialStringTok{f\textquotesingle{}Shape cinza: }\SpecialCharTok{\{}\NormalTok{img\_gray}\SpecialCharTok{.}\NormalTok{shape}\SpecialCharTok{\}}\SpecialStringTok{\textquotesingle{}}\NormalTok{)}

\NormalTok{mm.show([img, img\_gray], title}\OperatorTok{=}\NormalTok{[}\StringTok{\textquotesingle{}Original (RGB)\textquotesingle{}}\NormalTok{, }\StringTok{\textquotesingle{}Tons de Cinza\textquotesingle{}}\NormalTok{])}
\end{Highlighting}
\end{Shaded}

\subsection{1.4.2 Limiarização
(Binarização)}\label{limiarizauxe7uxe3o-binarizauxe7uxe3o-1}

A limiarização converte uma imagem em tons de cinza para binária:

\begin{equation}\phantomsection\label{eq-1-limiar}{g(x,y) = \begin{cases} 255 & \text{se } f(x,y) \geq T \\ 0 & \text{caso contrário} \end{cases}}\end{equation}

O método de \textbf{Otsu} \autocite{otsu1979} determina \(T\)
automaticamente maximizando a variância inter-classes do histograma.

\begin{Shaded}
\begin{Highlighting}[]
\CommentTok{\# Limiarização com T fixo}
\NormalTok{T }\OperatorTok{=} \DecValTok{128}
\NormalTok{img\_bin }\OperatorTok{=}\NormalTok{ mm.threshold(img\_gray, T)}

\CommentTok{\# Limiarização Otsu (T automático)}
\NormalTok{\_, img\_otsu }\OperatorTok{=}\NormalTok{ cv2.threshold(img\_gray, }\DecValTok{0}\NormalTok{, }\DecValTok{255}\NormalTok{,}
\NormalTok{                             cv2.THRESH\_BINARY }\OperatorTok{+}\NormalTok{ cv2.THRESH\_OTSU)}

\BuiltInTok{print}\NormalTok{(}\SpecialStringTok{f\textquotesingle{}Valores únicos (T=}\SpecialCharTok{\{}\NormalTok{T}\SpecialCharTok{\}}\SpecialStringTok{): }\SpecialCharTok{\{}\BuiltInTok{sorted}\NormalTok{(}\BuiltInTok{set}\NormalTok{(img\_bin.flatten()))}\SpecialCharTok{\}}\SpecialStringTok{\textquotesingle{}}\NormalTok{)}

\NormalTok{mm.show([img\_gray, img\_bin, img\_otsu],}
\NormalTok{        title}\OperatorTok{=}\NormalTok{[}\SpecialStringTok{f\textquotesingle{}Cinza\textquotesingle{}}\NormalTok{, }\SpecialStringTok{f\textquotesingle{}Binária (T=}\SpecialCharTok{\{}\NormalTok{T}\SpecialCharTok{\}}\SpecialStringTok{)\textquotesingle{}}\NormalTok{, }\StringTok{\textquotesingle{}Otsu\textquotesingle{}}\NormalTok{])}
\end{Highlighting}
\end{Shaded}

\section{1.5 Acesso a Pixels}\label{acesso-a-pixels-2}

Em Python/NumPy, o acesso a pixels usa indexação de arrays
\texttt{img{[}linha,\ coluna{]}}.

\begin{Shaded}
\begin{Highlighting}[]
\NormalTok{r, c }\OperatorTok{=} \DecValTok{100}\NormalTok{, }\DecValTok{100}

\CommentTok{\# Acessar pixel}
\NormalTok{pixel\_cinza }\OperatorTok{=}\NormalTok{ img\_gray[r, c]}
\NormalTok{pixel\_rgb   }\OperatorTok{=}\NormalTok{ img[r, c]       }\CommentTok{\# array [R, G, B]}

\BuiltInTok{print}\NormalTok{(}\SpecialStringTok{f\textquotesingle{}Pixel cinza (}\SpecialCharTok{\{}\NormalTok{r}\SpecialCharTok{\}}\SpecialStringTok{,}\SpecialCharTok{\{}\NormalTok{c}\SpecialCharTok{\}}\SpecialStringTok{): }\SpecialCharTok{\{}\NormalTok{pixel\_cinza}\SpecialCharTok{\}}\SpecialStringTok{\textquotesingle{}}\NormalTok{)}
\BuiltInTok{print}\NormalTok{(}\SpecialStringTok{f\textquotesingle{}Pixel RGB   (}\SpecialCharTok{\{}\NormalTok{r}\SpecialCharTok{\}}\SpecialStringTok{,}\SpecialCharTok{\{}\NormalTok{c}\SpecialCharTok{\}}\SpecialStringTok{): R=}\SpecialCharTok{\{}\NormalTok{pixel\_rgb[}\DecValTok{0}\NormalTok{]}\SpecialCharTok{\}}\SpecialStringTok{  G=}\SpecialCharTok{\{}\NormalTok{pixel\_rgb[}\DecValTok{1}\NormalTok{]}\SpecialCharTok{\}}\SpecialStringTok{  B=}\SpecialCharTok{\{}\NormalTok{pixel\_rgb[}\DecValTok{2}\NormalTok{]}\SpecialCharTok{\}}\SpecialStringTok{\textquotesingle{}}\NormalTok{)}

\CommentTok{\# Imagem sintética 5×5 com pixel central branco}
\NormalTok{syn }\OperatorTok{=}\NormalTok{ np.zeros((}\DecValTok{5}\NormalTok{, }\DecValTok{5}\NormalTok{), dtype}\OperatorTok{=}\NormalTok{np.uint8)}
\NormalTok{syn[}\DecValTok{2}\NormalTok{, }\DecValTok{2}\NormalTok{] }\OperatorTok{=} \DecValTok{255}
\BuiltInTok{print}\NormalTok{(}\StringTok{\textquotesingle{}}\CharTok{\textbackslash{}n}\StringTok{Matriz 5×5:\textquotesingle{}}\NormalTok{)}
\BuiltInTok{print}\NormalTok{(syn)}
\end{Highlighting}
\end{Shaded}

\section{1.6 Resumo}\label{resumo-2}

Neste capítulo foram apresentados os fundamentos de representação de
imagens digitais:

\begin{itemize}
\tightlist
\item
  Imagem digital = função \(f(x,y)\) mapeando coordenadas para
  intensidades
\item
  Tipos: binária (\(k=1\), \(\{0,255\}\)), tons de cinza (\(k=1\),
  \([0,255]\)), RGB (\(k=3\))
\item
  \texttt{morph.py}: biblioteca didática que simplifica operações de PDI
  em Python
\item
  Limiarização: operação fundamental de binarização
\end{itemize}

O Capítulo 2 estuda \textbf{histogramas e equalização}.

\subsection{Questões de Revisão}\label{questuxf5es-de-revisuxe3o-1}

\begin{enumerate}
\def\labelenumi{\arabic{enumi}.}
\tightlist
\item
  Qual a diferença entre imagem em tons de cinza e imagem binária?
\item
  Como o método de Otsu determina o limiar \(T\) automaticamente?
\item
  O que representa \texttt{img.shape} para uma imagem NumPy RGB?
\end{enumerate}

\printbibliography

\cleardoublepage
\phantomsection
\addcontentsline{toc}{part}{Appendices}
\appendix

\chapter*{Referências}\label{referuxeancias-1}
\addcontentsline{toc}{chapter}{Referências}

\markboth{Referências}{Referências}

\printbibliography[heading=none]


\backmatter



\end{document}
